\documentclass[11pt]{article}

\usepackage[a4paper,margin=29mm]{geometry}
\usepackage[T1]{fontenc}
\usepackage{lmodern}
\usepackage{microtype}
\usepackage{amsmath,amssymb,amsthm,mathtools}
\usepackage{array}
\usepackage{booktabs}
\usepackage{enumitem}
\usepackage{needspace}
\usepackage{float}
\usepackage{tikz}
\usetikzlibrary{arrows.meta,positioning}
\usepackage[hidelinks]{hyperref}

\newtheorem{theorem}{Theorem}[section]
\newtheorem{proposition}[theorem]{Proposition}
\newtheorem{corollary}[theorem]{Corollary}
\newtheorem{lemma}[theorem]{Lemma}
\theoremstyle{definition}
\newtheorem{definition}[theorem]{Definition}
\theoremstyle{remark}
\newtheorem{remark}[theorem]{Remark}

\newcommand{\F}{\mathbb{F}}
\newcommand{\cC}{\mathcal{C}}

\title{Quadratic trade rigidity and cubic orientation\\
in conic matching quotients}
\author{Tavis Rudd}
\date{July 2026}

\begin{document}
\begingroup
\makeatletter
\let\clebschmaintitle\@title
\def\@title{\large\textsc{Clebsch: Rigidity from Sparse Shadows --- }\textbf{II}\\[0.7em]
  \LARGE\clebschmaintitle}
\makeatother
\maketitle
\endgroup

\begin{abstract}
Restriction to a conic forgets how its marked points were paired into
secants.  Among full \(\operatorname{PGL}_2(q)\)-orbits of perfect matchings
over odd finite fields, we classify those whose conic-quotient evaluation
space has a one-dimensional strength-two trade---a signed relation
annihilating all quadratic coordinate products---generated by a two-valued
vector.  Exactly two occur: the balanced \(B_3/\F_7\) and \(H_3/\F_{11}\)
orbits.  The trade recovers their two complementary sheets, without assuming
self-association or Gorensteinness.

The restriction to perfect matchings is sharp.  For either surviving stabilizer, its fixed
locus in the ambient conic-product fiber is an affine line of \(q\) pairwise
nonconjugate rational points.  Only the matching point splits completely
into linear factors, although \(q-2\) nonmatching orbits have the same
one-dimensional two-valued strength-two trade condition; the remaining
rational point is the coalescence parameter.  Thus complete
splitting selects the matching orbit from the fixed line.  The exceptional
one-factorizations are classical; the new boundary is the fixed line and its
unique completely split point.

For the two matching configurations, the first nonzero signed tensor moment
is an anti-invariant cubic.  Their \(14\)- and \(22\)-point homogenizations
are self-associated and arithmetically Gorenstein, and maximal isotropy
identifies the cubic with the Macaulay inverse system of an Artinian
reduction.  General self-dual-code criteria already explain the Gorenstein
consequence of the Schur square; the configuration-specific contributions
are the all-field orbit classification, sharp matching boundary, sheet
reconstruction, and cubic orientation.  Targeted modular detectors, made
exhaustive by Faber's tame subgroup theorem, exclude every other matching
orbit without a field census.  A common alternating-cycle and Dickson
calculation supplies the radial nonvanishing needed for the quotient ranks.
\end{abstract}

\noindent\textbf{Keywords.} conic ideal; secant matching; strength-two
trade; one-factorization; Dickson polynomial; self-associated points;
arithmetically Gorenstein; cubic inverse system.

\smallskip
\noindent\textbf{2020 Mathematics Subject Classification.} 05B25, 51E22,
20C20, 13A02.

\section{Introduction}

A perfect matching of marked conic points determines a product of secant
lines.  Restriction to the conic deliberately forgets the pairing, and we ask
when low-degree algebra in the resulting quotient recovers it.  Over all odd
finite fields, a one-dimensional two-valued quadratic trade is rigid enough:
within the full matching carrier it forces exactly the $B_3/\F_7$ and
$H_3/\F_{11}$ orbits, and the first surviving odd tensor restores the
orientation of the recovered sheets.

\Needspace{27\baselineskip}
\medskip
\noindent\textbf{Reconstruction perspective.}
The classification ranges over every odd prime power.  It determines when the
affine quotient configuration recovers the missing pairing and when its first
odd tensor orients the recovered sheets.  The Clebsch-related
$H_3/\F_{11}$ case is one exceptional output, not an input.  The classification
and reconstruction are standalone; no companion construction is used here.

A four-endpoint switch shows that competing secant products differ by the
conic equation, so their quotients form an affine configuration.
For the two surviving orbits, quadratic products recover an unordered
bipartition of that configuration.  The first surviving signed tensor is
cubic and orients its two parts.  Thus quadratic data recover the
factorization without choosing an orientation; the first odd covariant
restores the choice.  By the \emph{matching carrier} we mean the locus of
products of the secant lines in a perfect matching of the marked conic
points.  This is a low-degree recognition theorem rather than an orbit census:
within the full matching carrier, the quadratic relation forces one of two
exceptional geometries without assuming its group type, self-association, or
Gorenstein coordinate ring.

The general implication from self-dual codes and their Schur squares to
Gorenstein coordinate rings is background here.  The theorem instead starts
from the weaker quadratic-trade condition, with no self-duality or Gorenstein
premise.

\begin{center}
\small
\begin{tabular}{@{}c@{\qquad}c@{\qquad}c@{}}
\toprule
type & linear quotient & sheets/cubic\\
\midrule
\(A_3/\F_5\) & rank \(3\) & not applicable\\
\(B_3/\F_7\) & rank \(6\) & recovered/oriented\\
\(H_3/\F_{11}\) & rank \(10\) & recovered/oriented\\
\bottomrule
\end{tabular}
\end{center}
\noindent
The all-field bottleneck is \emph{outer parity}.  A selected simple
\(H=\operatorname{PSL}_2(q)\)-detector has two
\(G=\operatorname{PGL}_2(q)\)-extensions, exchanged by the determinant
square-class character; the trade forces the opposite one into quadratic
moments.  Appendix A rules it out, Faber makes the detector cases exhaustive,
and Dickson enters only after the sheet size is known.

\begin{theorem}[Quadratic-trade classification and cubic orientation]
\label{thm:factorization-recovery}
Let the $A_3$, $B_3$, and $H_3$ matching configurations be the finite
matching orbits defined in Section~\ref{sec:rank-three}.  Then:
\begin{enumerate}[label=\textup{(\roman*)},leftmargin=*]
\item among all odd prime powers \(q\), the \(B_3/\F_7\) and
\(H_3/\F_{11}\) configurations are the only full perfect-matching
orbits under \(\operatorname{PGL}_2(q)\) whose
strength-two trade space is
one-dimensional and generated by a vector with exactly two values;
\item matching products scaled as in Section~\ref{sec:matching-products}
have equal restriction to the conic, and their conic-ideal quotients have
difference ranks $3$, $6$, and $10$;
\item in the $B_3$ and $H_3$ cases, the two factorization sheets are,
up to interchange, the unique complementary halves with equal first and
second tensor moments;
\item in those two cases, the signed moments vanish through degree two,
while a nonzero degree-three tensor transforms by the character that
exchanges the sheets;
\item the homogenized \(B_3\) and \(H_3\) configurations are
self-associated with Gorenstein coordinate rings; the signed cubic is
the Macaulay inverse system of their Artinian reductions;
\item for either surviving stabilizer, the fixed locus in the ambient
conic fiber is an affine line whose unique secant-product point is the
matching point, while \(q-2\) nonmatching orbits on that line still have
one-dimensional two-valued quadratic-trade space.
\end{enumerate}
\end{theorem}

Rodr\'iguez-Pajares, Ruano, and Salizzoni characterize the Gorenstein defect
of a self-dual code by its indecomposable blocks, equivalently by the defect of
its Schur square
\cite[Theorem~3.11 and Corollaries~3.12--3.13]{RodriguezPajaresEtAl2025}.
That general mechanism is not a priority claim here.  The new input and
outputs are the all-field trade classification, sheet reconstruction, sharp
matching boundary, and orbit-specific cubic orientation in the theorem above.

\Needspace{24\baselineskip}
\medskip
\noindent\textbf{Proof and reading map.}
\begin{center}
\small
\setlength{\tabcolsep}{3pt}
\renewcommand{\arraystretch}{0.9}
\begin{tabular}{@{}c c c@{}}
&
\fbox{\parbox{0.30\textwidth}{\centering
Matching products have a common conic restriction; differences define the
conic-ideal quotient.\\[-1pt]
\emph{Section~\ref{sec:matching-products}}}}
&\\
\(\swarrow\) && \(\searrow\)\\
\fbox{\parbox{0.25\textwidth}{\centering
Compute the three quotient ranks \(3,6,10\).\\[-1pt]
\emph{Section~\ref{sec:rank-three}; (ii)}}}
&&
\fbox{\parbox{0.29\textwidth}{\centering
A two-valued quadratic trade forces two
\(\operatorname{PSL}_2(q)\)-sheets.\\[-1pt]
\emph{Section~\ref{sec:rank-three}}}}
\\
\(\substack{\Downarrow\\[-2pt]\text{\scriptsize rank data used below}}\)
&& \(\Downarrow\)\\
& &
\fbox{\parbox{0.29\textwidth}{\centering
Faber's subgroup split is exhaustive; the detectors force each sheet to
have size \(q\).\\[-1pt]
\emph{Section~\ref{sec:rank-three}; Appendix~\ref{app:modular-detectors}}}}
\\
\(\vdots\) && \(\Downarrow\)\\
& &
\fbox{\parbox{0.29\textwidth}{\centering
Dickson's classification and the invariant-matching check identify the only
possible survivors as \(B_3/\F_7\) and \(H_3/\F_{11}\).\\[-1pt]
\emph{Section~\ref{sec:rank-three}; (i), necessity}}}
\\
\(\vdots\) &&\\[-2pt]
\(\searrow\) && \(\swarrow\)\\
&
\fbox{\parbox{0.32\textwidth}{\centering
The moment identities prove the converse in (i), recover the unordered
sheets, and show that the signed cubic orients them and gives Gorenstein
duality.\\[-1pt]
\emph{Section~\ref{sec:balanced}; (i), converse; (iii)--(v)}}}
&\\
& \(\Downarrow\) &\\
&
\fbox{\parbox{0.32\textwidth}{\centering
The fixed affine line exposes the off-carrier ambiguity; its inherited trade
and geometric splitting give the sharp carrier boundary.\\[-1pt]
\emph{Section~\ref{sec:balanced}; (vi)}}}
&
\end{tabular}
\end{center}
\noindent
The right-hand chain proves the necessity in (i); the two branches merge for
its converse and the recovery/orientation results.  The boundary inherits its
trade clause from that merged node.  Appendix A is required only at the
detector node.  Appendices B--F may be skipped on a first pass; Appendix G
records the trust boundary and supplies no logical premise.

For each surviving stabilizer, Section~\ref{sec:balanced} constructs a fixed
affine line inside the ambient conic fiber.  The forms in that ambient fiber
that split geometrically into linear factors comprise its geometric Chow
locus.  The carrier boundary is summarized by
\[
\begin{array}{rcl}
 \text{matching carrier and two-valued trade}
   &\Longrightarrow& B_3/\F_7\ \text{or}\ H_3/\F_{11},\\[2pt]
 \text{fixed conic-fiber line away from coalescence}
   &\Longrightarrow& q-1\ \text{trade-condition orbits},\\[2pt]
 \text{fixed conic-fiber line}\cap\text{geometric Chow locus}
   &=& \{\text{matching point}\}.
\end{array}                                                   \tag{1.1}
\]
The middle line consists of one matching orbit and \(q-2\) nonmatching
orbits; the excluded rational point is the coalescence parameter.  The last
line shows that imposing the Chow condition restores uniqueness.

\section{Conic matching products}
\label{sec:matching-products}

Write the standard conic as
\[
 \cC=V(Q),\qquad Q=XZ-Y^2,
\]
and parametrize it by $\nu(s:t)=(s^2:st:t^2)$.  For points
$a=(a_0:a_1)$ and $b=(b_0:b_1)$ of $\mathbb P^1$, put
$[a,b]=a_0b_1-a_1b_0$.  We scale the secant $L_{ab}$ through
$\nu(a)$ and $\nu(b)$ by the condition
\[
 L_{ab}(\nu(s:t))=[a,(s:t)]\,[b,(s:t)].
\]

Let $\Omega$ be a set of $2m$ distinct marked points of $\mathbb P^1$,
with fixed vector representatives.  For a perfect matching $M$ of
$\Omega$, define its matching product
\[
 P_M=\prod_{\{a,b\}\in M}L_{ab}.
\]
The vector representatives are part of the marked-conic data.  Their
role is to make the products comparable before passing to projective
equivalence.

This choice does not add projective structure.  If the representative of
each endpoint \(a\) is rescaled by \(\lambda_a\), then every matching product
is multiplied by the same scalar \(\prod_{a\in\Omega}\lambda_a\).  Hence every
quotient difference in (2.1) is multiplied by that common scalar: the affine
configuration changes by a homothety, while its trade data, ranks, and cubic
line are unchanged projectively.

\begin{proposition}[The matching-secant quotient]
\label{prop:matching-secant-quotient}
For every $m\ge2$ and every two perfect matchings $M,N$ of $\Omega$,
\[
 P_M|_{\cC}=P_N|_{\cC},\qquad P_M-P_N\in(Q).
\]
Consequently, after choosing a reference matching $M_0$,
\[
 \Phi_{M_0}(M)=\frac{P_M-P_{M_0}}{Q}
       \in H^0(\mathbb P^2,\mathcal O(m-2))
                                                        \tag{2.1}
\]
is an affine configuration.  Changing the reference matching translates
this configuration.  If $g\in\mathrm{GL}_2$ preserves $\Omega$,
$\rho(g)$ is its action on binary quadratics, and representatives are
transported by $g$, then
\[
 \Phi_{gM_0}(gM)\circ\rho(g)
   =\det(g)^{\,2m-2}\Phi_{M_0}(M).                       \tag{2.2}
\]
Thus the difference space, its rank, and its affine moment conditions do
not depend on $M_0$ and are equivariant under the endpoint stabilizer.
\end{proposition}

\begin{proof}
For $x\in\mathbb P^1$,
\[
 P_M(\nu(x))=\prod_{a\in\Omega}[a,x],
\]
which does not depend on the pairing.  The kernel of the Veronese
pullback
\[
 k[X,Y,Z]\longrightarrow k[s,t],\qquad
 (X,Y,Z)\longmapsto(s^2,st,t^2)
\]
is the principal ideal $(Q)$, so $Q$ divides every difference.
Changing the reference follows by subtraction.  Finally,
$L_{ga,gb}\circ\rho(g)=\det(g)^2L_{ab}$ and
$Q\circ\rho(g)=\det(g)^2Q$, which gives (2.2).
\end{proof}

\subsection{The four-endpoint switch}
\label{sec:switches}

For four distinct endpoints, the Pl\"ucker relation gives the local
switch identity
\[
 L_{ab}L_{cd}-L_{ac}L_{bd}
   =[a,d][b,c]\,Q.                                      \tag{2.3}
\]
Multiplying (2.3) by the secant factors common to two matchings proves
divisibility for a single switch.  Any two perfect matchings are joined
by a sequence of such switches: if $M$ and $N$ differ at $a$, write
$\{a,b\}\in M$ and $\{a,c\}\in N$; switching the two $M$-edges through
$b$ and $c$ creates $\{a,c\}$ and reduces the symmetric difference.
This also gives an elementary proof of the divisibility clause in
Proposition~\ref{prop:matching-secant-quotient}.

The switch calculation and the global pullback proof play different
roles.  The pullback proves the quotient for arbitrary matchings at once.
The switch identity records how local changes of factorization move in
the quotient space.  A selected Coxeter orbit need not contain every
perfect matching or every intermediate switch.

\section{Matching-orbit classification and quotient ranks}
\label{sec:rank-three}

The three distinguished finite matching orbits are discrete input:
Proposition~\ref{prop:matching-secant-quotient} applies to every matching,
but it does not distinguish the Coxeter orbits.
Here and below, ``rank three'' refers to Coxeter rank.  It does not
refer to the dimensions of the conic-ideal quotient spaces, which are
\(3,6,10\).

For $q\in\{5,7,11\}$, index
\[
 \mathbb P^1(\F_q)=\{0,1,\ldots,q-1,\infty\},
\]
where $a$ denotes $(1:a)$ and $\infty$ denotes $(0:1)$.  Let
$G_q=\operatorname{PGL}_2(q)$ act in the usual way.  The three base
matchings are
\[
\begin{array}{c@{\qquad}c@{\qquad}l}
T&q&M_T\\ \hline
A_3&5&
 \{0,\infty\},\{1,4\},\{2,3\},\\
B_3&7&
 \{0,2\},\{1,4\},\{3,\infty\},\{5,6\},\\
H_3&11&
 \{0,1\},\{2,5\},\{3,7\},\{4,9\},\{6,8\},
 \{10,\infty\}.
\end{array}                                               \tag{3.1}
\]
Define the \emph{rank-three matching configuration}
\[
 \Omega_T=G_q\cdot M_T.
\]
The stabilizers of the three displayed matchings are respectively
$S_4,S_4,A_5$.  For a direct check, with
\(\left(\begin{smallmatrix}a&b\\c&d\end{smallmatrix}\right)\) acting by
\(x\mapsto(ax+b)/(cx+d)\), generators preserving the three rows of (3.1)
may be chosen as
\[
\begin{array}{c|c|c|c}
T&a&b&\text{relations}\\
\hline
A_3&\left(\begin{smallmatrix}0&1\\2&0\end{smallmatrix}\right)
&\left(\begin{smallmatrix}1&1\\2&3\end{smallmatrix}\right)
&a^2=b^3=(ab)^4=1\\[3pt]
B_3&\left(\begin{smallmatrix}0&1\\3&2\end{smallmatrix}\right)
&\left(\begin{smallmatrix}1&1\\1&5\end{smallmatrix}\right)
&a^3=b^4=(ab)^2=1\\[3pt]
H_3&\left(\begin{smallmatrix}0&1\\2&3\end{smallmatrix}\right)
&\left(\begin{smallmatrix}0&1\\7&8\end{smallmatrix}\right)
&a^3=b^5=(ab)^2=1.
\end{array}
\]
These are the standard triangle presentations of \(S_4,S_4,A_5\);
the maximal-subgroup classification and the outer-normalizer check
\cite{Giudici2007} show that they are the full stabilizers in \(G_q\).
Thus orbit--stabilizer gives
\[
 |\Omega_{A_3}|=5,\qquad |\Omega_{B_3}|=14,\qquad
 |\Omega_{H_3}|=22.                                      \tag{3.2}
\]
These matching orbits and their stabilizers are classical.  In particular,
Edge and Dye identify the exceptional finite conic configurations and
substantial parts of their incidence geometry
\cite{Edge1956,Dye1991}.  The exceptional one-factorizations themselves
also precede this work: Cameron and Korchm\'aros classify the doubly
transitive case, and Han's classification of symmetric factorizations
contains the triples \(\operatorname{PSL}_2(7),S_4,8,7\) and
\(\operatorname{PSL}_2(11),A_5,12,11\)
\cite{CameronKorchmaros1993,Han2025}.  We use the Coxeter names for the
corresponding tetrahedral, octahedral, and icosahedral data; neither the
orbit counts in (3.2) nor the ambiguity of an unmarked conic is asserted
here as new.

Two adjacent frameworks should also be separated from the claim here.
Korchm\'aros, Pace, and Sonnino construct one-factorizations from partitions
of the external points of an oval, each part meeting every tangent once
\cite{KorchmarosPaceSonnino2018}; the present carrier instead consists of
secant matchings on the marked conic, followed by passage through its ideal.
Likewise, the group--normal-subgroup architecture in which one group fixes
the factors and a larger group permutes them is the classical homogeneous
factorization framework of Giudici, Li, Poto\v{c}nik, and Praeger
\cite{GiudiciLiPotocnikPraeger2006}.  The assertions below concern the
specific conic-quotient and trade-rigidity mechanism, not that architecture.

\subsection{Trade reduction and the modular obstruction}

Here is the application dictionary for the abstract reduction.  We take
\(G=\operatorname{PGL}_2(q)\), \(H=\operatorname{PSL}_2(q)\), and
\(\Omega\) to be a full matching orbit; once the two trade levels are
identified with its \(H\)-orbits, they are the candidate sheets.  The module
\(\mathcal A\) is the space of ambient affine-linear functions,
\(E=\mathcal A^\vee\) is its homogenized point-vector module, and
\(F=\ker(E\to k)=\operatorname{Sym}^{(q-3)/2}L(2)\) is the translation
module.  The quadratic moment map is
\(\mu([M])=v_M^2\).  The next lemma converts the
trade relation into simultaneous occurrences of the two \(G\)-extensions of
one \(H\)-detector in \(\operatorname{Sym}^2E\); the uniform exclusion later
rules out the opposite-parity occurrence by eliminating both its zero-linear
and split-pullback alternatives.

\begin{lemma}[Projective--trade reduction]
\label{lem:projective-trade-reduction}
Let \(k\) be an algebraically closed field of odd characteristic, let
\(H\triangleleft G\) have index two, and let a transitive \(G\)-set
\(\Omega\) split as two transitive \(H\)-sets
\(\Omega_+\sqcup\Omega_-\).  Let
\[
 0\longrightarrow F\longrightarrow E
  \mathrel{\mathop{\longrightarrow}^{\epsilon}}k
  \longrightarrow0
\]
be an exact sequence of \(kG\)-modules, and let
\(v:\Omega\to E\) be a \(G\)-map with
\(\epsilon(v_\omega)=1\).  Define
\[
 \mu:k[\Omega]\longrightarrow\operatorname{Sym}^2E,
 \qquad [\omega]\longmapsto v_\omega^2.
\]
Also put
\[
 a:k[\Omega]\longrightarrow E,\qquad [\omega]\longmapsto v_\omega.
\]
All embeddings, pullbacks, and splittings below are in the category of
\(kG\)-modules.
Assume that
\[
 \ker\mu
 =k\left(\sum_{\omega\in\Omega_+}[\omega]
          -\sum_{\omega\in\Omega_-}[\omega]\right),
\]
and that \(k[\Omega_+]\) contains a nontrivial simple
\(H\)-submodule \(S\) such that \({}^gS\simeq S\) for
\(g\in G\smallsetminus H\), and that \(S\) extends to \(G\).
If \(S^+\) and \(S^-=S^+\otimes\chi\) are its two
\(G\)-extensions, where \(\chi\) is the nontrivial character of \(G/H\),
then both \(S^+\) and \(S^-\) embed in
\(\operatorname{Sym}^2E\).  For either summand \(T=S^\pm\), let
\(j=\mu|_T:T\to\operatorname{Sym}^2E\) be the resulting embedding and
put \(i=\partial j\).  Exactly one of the following occurs:
\begin{enumerate}
\item \(i=0\) and \(j(T)\subseteq\operatorname{Sym}^2F\);
\item \(i:T\to E\) is an embedding and the pullback along \(i\) of
  \[
   0\longrightarrow\operatorname{Sym}^2F
   \longrightarrow\operatorname{Sym}^2E
   \mathrel{\mathop{\longrightarrow}^{\partial}}E
   \longrightarrow0
  \]
  splits, where
  \[
   \partial(xy)=\epsilon(x)y+\epsilon(y)x.
  \]
\end{enumerate}
Moreover, let \(\xi\in\operatorname{Ext}^1_G(k,F)\) be the class of
the affine extension and let
\(\pi:F\otimes F\to\operatorname{Sym}^2F\) be the canonical quotient,
\(\pi(f\otimes f')=ff'\).  Since
\(T\) is nontrivial, \(i(T)\subseteq F\), and
\[
 \delta_T(i)
 =i^*\pi_*(F\otimes\xi)
 \quad\text{in}\quad
 \operatorname{Ext}^1_G(T,\operatorname{Sym}^2F).               \tag{3.2a}
\]
\end{lemma}

\begin{proof}
Choose \(g\in G\smallsetminus H\).  The \(G\)-submodule of
\(k[\Omega]\) generated by \(S\) is \(S\oplus{}^gS\).  It has zero
intersection with the displayed kernel, since that kernel is trivial
as an \(H\)-module and \(S\) is not.  Thus \(\mu\) embeds
\(S\oplus{}^gS\).

Choose the extension \(S^+\).  The elementary tensor identity for
induction gives
\[
 \operatorname{Ind}_H^G S
\mathrel{\mathop{\longrightarrow}^{\sim}}S\oplus{}^gS
 \simeq S^+\otimes k[G/H].
\]
The first map sends \(x\otimes s\) to \(xs\); it is an isomorphism
because its two sheet-supported images are disjoint.
Since \(2\ne0\) in \(k\), one has
\(k[G/H]\simeq k\oplus k_\chi\).  Hence the last module is
\(S^+\oplus S^-\), proving that both parities embed.

It remains only to inspect the canonical map \(\partial\).  It is
\(G\)-equivariant, and on every basis vector
\[
 \partial\mu([\omega])=2v_\omega=2a([\omega]).
\]
Choose \(e\in E\) with \(\epsilon(e)=1\), merely as
a vector-space splitting.  Then
\[
 \operatorname{Sym}^2E
 =ke^2\oplus eF\oplus\operatorname{Sym}^2F,
 \qquad
 \partial(e^2)=2e,\quad \partial(ef)=f.
\]
Consequently the displayed sequence is exact.  Since \(T\) is simple,
\(i\) is either zero or injective.  In the first case the image
of \(j\) lies in \(\operatorname{Sym}^2F\).  In the second, the map
\(s\mapsto(j(s),s)\) is a section of the pullback of that sequence along
\(i\).  Equivalently, if
\[
 \delta_T:\operatorname{Hom}_G(T,E)
 \longrightarrow\operatorname{Ext}^1_G
   (T,\operatorname{Sym}^2F)
\]
is the connecting map, then \(\delta_T(i)=0\).  This proves the
dichotomy.  We use the convention that \(\delta_T(i)\) is the Yoneda
class of the pullback along \(i\).

For the final assertion, put \(D=\partial^{-1}(F)\).  Multiplication
\(f\otimes x\mapsto fx\) gives a morphism from
\[
 0\longrightarrow F\otimes F\longrightarrow F\otimes E
 \mathrel{\mathop{\longrightarrow}^{1\otimes\epsilon}}F
 \longrightarrow0
\]
onto
\[
 0\longrightarrow\operatorname{Sym}^2F\longrightarrow D
 \mathrel{\mathop{\longrightarrow}^{\partial}}F
 \longrightarrow0.
\]
The left map is \(\pi\), the right map is the identity, and the lower
row is the pushout of the upper row: this is immediate from the
vector-space splitting \(E=ke\oplus F\) used above.  Thus the lower
row represents \(\pi_*(F\otimes\xi)\).  Since
\(\epsilon i:T\to k\) is zero for nontrivial simple \(T\), the map \(i\)
lands in \(F\); pulling back the lower row gives (3.2a).
More concretely, choose the same vector-space lift \(e\) and put
\(c(g)=ge-e\in F\).  With the connecting-cocycle convention
\(g\mapsto gsg^{-1}-s\), the class in (3.2a) is represented by
\[
 z_i(g):T\longrightarrow\operatorname{Sym}^2F,\qquad
 z_i(g)(t)=i(t)c(g).                                           \tag{3.2b}
\]
Changing \(e\) changes \(z_i\) by a coboundary.
\end{proof}

\subsection{Detector modules and their outer parity}

The all-field exclusion needs no classification of the full finite-group
socle of \(F\).  It uses only two extension-field linear channels, the
prime-field Fischer summands, and the absence of one outer parity from the
quadratic moment module.  We state exactly those interfaces here; their
specialist proofs are collected in
Appendix~\ref{app:modular-detectors}.

\begin{lemma}[Targeted linear detectors]
\label{lem:targeted-linear-detectors}
Let \(q=p^e\) be odd, let \(H=\operatorname{PSL}_2(q)\), and put
\[
 d=\frac{q-3}{2},\qquad F=\operatorname{Sym}^{d}L(2).
\]
\begin{enumerate}
\item For the Steinberg module \(\operatorname{St}=L(q-1)\),
\[
 \operatorname{Hom}_H(\operatorname{St},F)=0.
\]
\item If \(q\equiv3\pmod4\), \(e>1\), and \(S=L(q-7)\), then
\[
 S\lhook\joinrel\longrightarrow F.
\]
Every \(H\)-map in this channel is algebraic.  In the covariant
polynomial convention the displayed map is
\(L(q-7)\otimes\det^2\to F\); after normalizing the target by
\(\det^{-d}\), its source twist is
\(\det^{2-d}=\det^{-(q-7)/2}\).  Thus every occurrence has this
determinant-normalized outer parity, and the other
\(\operatorname{PGL}_2(q)\)-extension is absent from \(F\).
\item If \(e=1\) and \(p\geq5\), then
\[
 F=\bigoplus_{0\leq j\leq\lfloor(p-3)/4\rfloor}
       L(p-3-4j),
\]
with every summand occurring once and admitting an \(H\)-retraction.
\end{enumerate}
\end{lemma}

\begin{proof}
See Appendix~\ref{app:linear-detectors}.  The root-degree bound is
\(q-1\) in the two extension-field channels, so their finite maps are
algebraic.  The \(L(q-7)\) map is multiplication of a digitwise Cartan
embedding by the discriminant, and the prime-field statement follows from
the tilting Fischer character identity.
\end{proof}

\begin{lemma}[Opposite-parity quadratic obstruction]
\label{lem:opposite-parity-quadratic}
Let \(S\) be one of the detector modules used below: the Steinberg module,
\(L(q-7)\) in the extension-field branch, or one of the prime-field modules
\(L(s)\), \(s\in\{2,4,6,8,12\}\), in the range displayed in (3.2f).
Let \(S^\square\) be the extension to
\(G=\operatorname{PGL}_2(q)\) opposite to the covariant
determinant-normalized extension.  Then
\[
 \operatorname{Hom}_G(S^\square,\operatorname{Sym}^2F)=0,
\]
and the other extension is \(S^\square\otimes\chi\), where \(\chi\) is
the nontrivial character of \(G/H\).
\end{lemma}

\begin{proof}
Appendix~\ref{app:outer-parity-proof} constructs an injection from this
finite-group Hom space into
\[
 \operatorname{Hom}_{\operatorname{SL}_2(k)}
 \bigl(S\otimes L(1)\otimes L(1)^{(e)},\operatorname{Sym}^2F\bigr)
\]
and proves the latter space zero for each listed source.  The Steinberg
case uses the complete two-section Weyl filtration of \(T(2q)\), not only
its head.
\end{proof}

\begin{lemma}[Affine-class contraction]
\label{lem:affine-class-contraction}
Let \(i:S^\square\hookrightarrow F\) have an \(H\)-retraction
\(\rho:F\to S\).  For the point-vector extension
\(0\to F\to E\to k\to0\), with class \([c]\), the pullback along \(i\)
of
\[
 0\longrightarrow\operatorname{Sym}^2F
 \longrightarrow\operatorname{Sym}^2E
 \mathrel{\mathop{\longrightarrow}^{\partial}}E
 \longrightarrow0
\]
contracts to
\[
 (\dim S)[c]+i_*\rho_*[c]\in H^1(H,F).                 \tag{3.2c}
\]
For the prime-field Fischer detectors, \([c]\ne0\), it is supported only
on the top summand \(L(p-3)\), and the resulting scalar in (3.2c) is
nonzero.
\end{lemma}

\begin{proof}
The evaluation--coevaluation calculation, the polynomial representative
of \([c]\), and the Borel-cohomology support calculation are given in
Appendix~\ref{app:contraction-proof}.
\end{proof}

\subsection{Uniform sheet exclusion}

\begin{lemma}[Uniform nonprincipal-sheet exclusion]
\label{lem:uniform-sheet-exclusion}
Let \(q=p^e\) be odd, let \(G=\operatorname{PGL}_2(q)\) and
\(G^+=\operatorname{PSL}_2(q)\), and let \(\Omega\) be a full
\(G\)-orbit of perfect
matchings split into two \(G^+\)-orbits.  Let
\(L\subseteq\F_q^\Omega\) be its affine evaluation space.  If
\((L^{\circ2})^\perp\subseteq\F_q^\Omega\) is the sheet-sign line,
then either sheet has \(q\) members.
\end{lemma}

\begin{proof}
We give the modular argument because it removes the one-factorization
hypothesis.  Let \(K\leq G^+\) stabilize a matching and write a sheet as
\(G^+/K\).  A nontrivial unipotent cannot fix a perfect matching: it has
one fixed endpoint, while the partner of that endpoint would have an
orbit of length \(p\).  Thus \(K\) is a \(p'\)-group and
\[
 P_0=\F_q[G^+/K]
\]
is projective.  Its restriction to a translation Sylow subgroup is a
sum of regular modules, and
\[
 |G^+/K|=q\lambda                                      \tag{3.2c$_0$}
\]
for some positive integer \(\lambda\).  Extend scalars to
\(k=\overline{\F}_q\) and put \(P=k\otimes_{\F_q}P_0\).

Let \(\mathcal A\) be the full space of affine-linear functions on the
ambient conic-quotient form space, after scalar extension to \(k\), and let
\(\mathcal A\twoheadrightarrow L_k\) be restriction to \(\Omega\).  Put
\(E=\mathcal A^*\).  Thus \(E\) is the ambient point-vector module; it is
not the dual of the possibly smaller orbit span \(L_k\).  The inclusion of
the constants in \(\mathcal A\) dualizes to
\[
 0\longrightarrow F\longrightarrow E
 \mathrel{\mathop{\longrightarrow}^{\epsilon}}k
 \longrightarrow0,\qquad
 F=\operatorname{Sym}^{(q-3)/2}L(2),
                                                               \tag{3.2d}
\]
and evaluation at a matching \(M\) is a vector \(v_M\in E\) with
\(\epsilon(v_M)=1\).  The evaluation of products of functions in
\(\mathcal A\) factors as
\[
 \operatorname{Sym}^2\mathcal A\twoheadrightarrow
 \operatorname{Sym}^2L_k\longrightarrow k^\Omega.
\]
Its transpose is the moment map in
Lemma~\ref{lem:projective-trade-reduction}.  Since the first arrow is
surjective, the kernel of that transpose is exactly the annihilator of
\(L_k^{\circ2}\), hence the scalar extension of the strength-two trade
space.  Here we use that scalar extension is flat and commutes with
\(\operatorname{Sym}^2\).  This ambient formulation is essential when the
orbit spans a proper subspace, as it does for \(H_3\).
Write \(\mu\) for this moment map and
\(a:k[\Omega]\to E\) for the linear moment map
\([M]\mapsto v_M\).

We apply Lemmas~\ref{lem:targeted-linear-detectors},
\ref{lem:opposite-parity-quadratic}, and
\ref{lem:affine-class-contraction}.  Denote the absent outer extension by
\(S^\square\), and write
\[
 j_\square=\mu|_{S^\square},\qquad
 i_\square=\partial j_\square=2a|_{S^\square}.
\]
If \(i_\square=0\), then \(j_\square\) would embed \(S^\square\) in
\(\operatorname{Sym}^2F\), contradicting
Lemma~\ref{lem:opposite-parity-quadratic}.  We choose \(S\) so that this
already occurs over extension fields.  Over prime fields, a selected
linear copy instead retracts and Lemma~\ref{lem:affine-class-contraction}
rules out the split pullback.

The list of possible stabilizers is exhaustive without a recursive
subfield argument.  Since \(K\) is \(p'\), it is a finite \(p\)-regular
subgroup of \(\operatorname{PGL}_2(\overline{\F}_p)\).  Faber's tame
subgroup theorem \cite[Theorem~C]{Faber2023} gives
\[
 K\text{ cyclic, dihedral, }A_4,S_4,\text{ or }A_5.     \tag{3.2e}
\]
Only this abstract type is used in the following detector split.

Suppose first that \(e>1\).  If a cyclic or dihedral \(K\) is
intransitive on \(\mathbb P^1(\F_q)\), take
\(S=\operatorname{St}=L(q-1)\).  The endpoint permutation module is
\[
 k[\mathbb P^1(\F_q)]=k\oplus\operatorname{St},
\]
so \(\dim S^K=|K\backslash\mathbb P^1(\F_q)|-1>0\).
Lemma~\ref{lem:targeted-linear-detectors} gives
\(\operatorname{Hom}_H(S,F)=0\); hence \(i_\square=0\), a contradiction.

It remains among cyclic and dihedral groups only to consider a transitive
dihedral \(K\).  A cyclic subgroup of \(H\) lies in a nonsplit torus of
order at most \((q+1)/2\), so it cannot be transitive on \(q+1\)
endpoints.  Thus \(K\) is the full nonsplit normalizer of order \(q+1\),
acting regularly.  Put \(n=(q+1)/2\), choose a generator \(z\) of the
nonsplit torus of order \(2n\) in \(G\), with \(z\notin G^+\), and write
\[
 r=z^2,\qquad
 K=\langle r,w:r^n=w^2=1,\ wrw=r^{-1}\rangle .
\]
A \(K\)-invariant matching on the regular \(K\)-set is right multiplication
by a nonidentity involution.  Since \(zwz^{-1}=rw\), conjugation by
\(r^az\) sends \(r^jw\) to \(r^{j+2a+1}w\), preserving the matching
exactly when
\[
 2a+1=0\pmod n.                                         \tag{3.2f$_0$}
\]
If \(n\) is even, the rotation involution \(r^{n/2}\) is fixed by \(z\);
if \(n\) is odd, (3.2f$_0$) is soluble for every reflection matching.
Thus a split dihedral orbit can only be a reflection matching with \(n\)
even.  Hence \(q\equiv3\pmod4\), and we take \(S=L(q-7)\).
Its unique zero-weight line has Weyl sign \(+1\), so \(S^K\ne0\).
Lemma~\ref{lem:targeted-linear-detectors} puts every linear occurrence in
the determinant-normalized extension; therefore
\(a|_{S^\square}=0\), again a contradiction.

If \(K\) is exceptional and intransitive, the same Steinberg argument
applies:
\[
 \dim\operatorname{St}^K
 =|K\backslash\mathbb P^1(\F_q)|-1>0.
\]
If it is transitive, then \(q+1\mid |K|\), while
\(|K|\in\{12,24,60\}\).  For an odd nonprime prime power this leaves only
\(q=9\), where \(p=3\) divides every exceptional order, contrary to the
\(p'\) hypothesis.  This closes the extension-field cases.

Now let \(e=1\).  For cyclic, dihedral, and exceptional stabilizers use
\[
\begin{array}{c|c}
\text{type and range}&S\\ \hline
\text{cyclic},\ p\ge5&L(2)\\
\text{dihedral},\ p\ge7&L(4)\\
A_4,\ p\ge11&L(6)\\
S_4,\ p\ge11&L(8)\\
A_5,\ p\ge17&L(12).
\end{array}                                                \tag{3.2f}
\]
The cyclic and dihedral modules have the required zero-weight line, with
Weyl signs \(-1\) and \(+1\).  For the exceptional groups, averaging the
ordinary symmetric-power character over the binary-polyhedral lift gives
\[
\begin{array}{c|c|c}
K&S&|K|\dim S^K\\ \hline
A_4&L(6)&12\\
S_4&L(8)&24\\
A_5&L(12)&60.
\end{array}
\]
Here the \((\text{class size},\text{lift order})\) pairs are respectively
\[
\begin{array}{c|l}
A_4&(1,1),(3,4),(8,6)\\
S_4&(1,1),(9,4),(8,6),(6,8)\\
A_5&(1,1),(15,4),(20,6),(12,10),(12,5).
\end{array}
\]
Thus \(\dim S^K=1\) in every exceptional row.  Since \(K\) is \(p'\),
averaging is valid in \(k\).  Frobenius reciprocity and the self-dual
permutation form embed each self-dual \(S\) in \(P=k[G^+/K]\).

By Lemma~\ref{lem:targeted-linear-detectors}, each selected \(S\) is either
absent from \(F\), giving \(i_\square=0\), or is a multiplicity-one Fischer
summand.  In the latter case \(i_\square=0\) gives the same immediate
contradiction; otherwise \(i_\square\) identifies the unique copy and its
Fischer projection is an \(H\)-retraction.
Lemma~\ref{lem:affine-class-contraction} then makes the pulled-back affine
class nonzero.  Lemma~\ref{lem:opposite-parity-quadratic} excludes the
alternative quadratic embedding in both cases.

The case tree is now summarized by the detector or obstruction used in
each branch:
\begin{center}
\small
\begin{tabular}{@{}c@{\quad}l@{\quad}l@{}}
\toprule
regime & stabilizer & detector or obstruction\\
\midrule
\(e>1\) & intransitive cyclic, dihedral, exceptional
  & \(\operatorname{St}\), zero linear moment\\
\(e>1\) & full nonsplit normalizer
  & \(L(q-7)\), opposite linear parity\\
\(e>1\) & transitive exceptional
  & only \(q=9\), but then \(p\mid |K|\)\\
\(e=1\) & cyclic, dihedral, \(A_4,S_4,A_5\)
  & (3.2f), then absence or contraction\\
small & \(q=3,5,7,13\)
  & one orbit, outer stabilizer, or no subgroup\\
\bottomrule
\end{tabular}
\end{center}
The remaining endpoint statements make the last row explicit:
\[
\begin{array}{c|c|c|c|c}
q&K&\text{orbits}&\text{unique matching data}&N_G(K)\\ \hline
5&A_4&6&C_2<V_4\text{ is the two-point block}&S_4\\
7&A_4&4+4&
 N_{A_4}(C_3)=C_3\text{ gives the cross-orbit pairing}&S_4\\
5&D_6&6&\text{reflection, }n=3\text{ in (3.2f$_0$)}&D_{12}.
\end{array}                                               \tag{3.2g$_0$}
\]
In each row the displayed outer normalizer preserves the unique
matching, so the full orbit does not give a split \(\lambda>1\) case.
For \(q=3\), the
three perfect matchings form one \(G^+\)-orbit, so the two-sheet premise
is impossible.  There is no \(A_5\) row at \(q=13\), since
\(60\nmid|\operatorname{PSL}_2(13)|\).  We have therefore produced the
quadratic obstruction for every
possible split stabilizer with \(\lambda>1\), a contradiction.
Hence \(\lambda=1\).
\end{proof}

\subsection{Balanced-orbit completeness}

\begin{theorem}[Uniform balanced-orbit classification]
\label{thm:balanced-orbit-completeness}
Let \(q\) be an odd prime power, put
\[
 G=\operatorname{PGL}_2(q),\qquad G^+=\operatorname{PSL}_2(q),
\]
and let \(\Omega\) be a \(G\)-orbit of perfect matchings of
\(\mathbb P^1(\F_q)\).  Form its conic-ideal quotient configuration and
let \(L\subseteq\F_q^\Omega\) be its affine evaluation space.  Put
\[
 L^{\circ2}=\operatorname{span}\{fg:f,g\in L\},
\]
with pointwise multiplication, and take its annihilator under the
standard coordinate pairing on \(\F_q^\Omega\).  Suppose
\((L^{\circ2})^\perp\) is one-dimensional and every nonzero vector in it
takes exactly two values.  Then
\[
 (q,\Omega)=(7,\Omega_{B_3})
 \quad\text{or}\quad
 (q,\Omega)=(11,\Omega_{H_3}).
                                                        \tag{3.2h}
\]
Conversely, both displayed orbits satisfy these intrinsic hypotheses.
For either orbit, the two level sets are its two one-factorization sheets.
\end{theorem}

\begin{proof}
Proposition~\ref{prop:matching-secant-quotient} makes \(L\), its square,
and the line
\[
 R=(L^{\circ2})^\perp
\]
\(G\)-stable.  Choose \(0\ne\tau\in R\).  Its zero set is
\(G\)-invariant, so transitivity gives \(\tau\) full support.  The two
level sets of \(\tau\) depend only on \(R\), hence \(G\) permutes them.
This action is nontrivial because \(G\) is transitive on \(\Omega\).

For \(q>3\), the commutator subgroup of \(G\) is \(G^+\), and \(G/G^+\)
has order two.  The kernel of the action on the two levels is therefore
\(G^+\).  The stabilizer of a point fixes its level and so lies in
\(G^+\); consequently the two levels are exactly the two \(G^+\)-orbits.
An outer element acts on \(R\) by a scalar and exchanges the two nonzero
values of \(\tau\), so those values are opposites.

Lemma~\ref{lem:uniform-sheet-exclusion} shows that each intrinsic level
has size \(q\).  Thus \(|\Omega|=2q\).  When \(q=3\), there are only
three perfect matchings altogether, so such an orbit cannot occur.

Now fix \(M\in\Omega\) and let \(K=\operatorname{Stab}_G(M)\).
Orbit--stabilizer and the derived two-sheet conclusion give
\[
 |K|=\frac{q^2-1}{2},\qquad K\leq G^+,\qquad [G^+:K]=q.    \tag{3.2i}
\]

Assume \(q\geq5\), and place \(K\) in a maximal subgroup of \(G^+\).
Dickson's maximal-subgroup theorem, in the form recorded by Giudici
\cite[Theorem~2.2]{Giudici2007}, leaves Borel, dihedral, subfield, and
exceptional \(A_4,S_4,A_5\) subgroups.  The first two have order smaller
than \((q^2-1)/2\).  A proper subfield subgroup is no larger than
\(\operatorname{PGL}_2(r)\), where \(q=r^e\) and \(e\geq2\), and
\[
 r(r^2-1)<\frac{r^{2e}-1}{2}.
\]
Thus it too is too small.  An exceptional maximal subgroup has order at
most \(60\), so (3.2i) leaves \(q=5,7,9,11\).  For \(q=9\), the required
order \(40\) divides none of \(12,24,60\).  In the other three cases,
order and divisibility force
\[
 (q,K)=(5,A_4),\ (7,S_4),\ (11,A_5).                    \tag{3.2j}
\]

It remains to ask which of these subgroup actions preserves a matching.
This is a block-system question, not a matching census.  For \(q=5\),
the \(A_4\)-action on the six endpoints is transitive with point
stabilizer \(C_2\).  Its normal Klein four subgroup gives the unique
two-point block system.  The normalizer of \(A_4\) in
\(\operatorname{PGL}_2(5)\) is \(S_4\); uniqueness makes this whole
normalizer preserve the matching.  Its full orbit therefore has five
points, namely \(\Omega_{A_3}\), rather than ten.  In particular, the
other ten matchings cannot form a balanced \(5+5\) orbit.

For \(q=7\), the \(S_4\)-action on eight endpoints has point stabilizer
\(C_3\).  A two-point block stabilizer must have order six, and the
unique possibility containing that \(C_3\) is its normalizer \(S_3\).
Hence the invariant matching is unique.  Similarly, for \(q=11\), the
\(A_5\)-action on twelve endpoints has point stabilizer \(C_5\); its
unique possible two-point block stabilizer is
\(N_{A_5}(C_5)=D_{10}\).  The invariant matching is again unique.
In each case the two \(G^+\)-classes of the exceptional subgroup are
fused by \(G\), and the subgroup has no outer normalizer
\cite[Lemma~2.3]{Giudici2007}: the diagonal outer automorphism fuses the
two classes, so no element in its coset normalizes a representative.
Indeed, if \(J=\operatorname{Stab}_G(M)\), then
\(J\cap G^+\) contains \(K\).  Maximality gives
\(J\cap G^+=K\) or \(G^+\), and the latter is impossible because
\(G^+\) does not fix \(M\).  If \(J>K\), an element of \(J\setminus K\)
would therefore be an outer normalizer of \(K\), contrary to the cited
fusion statement.  Thus \(J=K\), so the orbit has size \(2q\) and
splits into the two claimed \(G^+\)-orbits.  The unique block systems
are precisely the matchings displayed in (3.1), up to
\(G\)-conjugacy.

Conversely, \(q=7,11\) are both \(3\pmod4\), so \(G^+\) is transitive on
unordered pairs of endpoints.  Each \(q\)-matching sheet therefore covers
every edge equally often; its total of \(q(q+1)/2\) edge occurrences
equals the number of edges, so it is a one-factorization.
Propositions~\ref{prop:radical-hadamard} and
\ref{prop:modular-sheet-mechanism} identify
\((L^{\circ2})^\perp\) with the sheet-sign line.  Hence both displayed
orbits satisfy the intrinsic hypotheses.
\end{proof}

Fix $M_T$ as reference and apply (2.1) to every member of $\Omega_T$.
For \(j\ge0\), write \(R_j=\F_q[X,Y,Z]_j\), and set \(R_j=0\) for
\(j<0\).  Put
\[
 \Delta_Q=4\partial_X\partial_Z-\partial_Y^2,\qquad
 \mathcal H_d=\ker(\Delta_Q:R_d\longrightarrow R_{d-2}).
                                                               \tag{3.3}
\]
The operator $\Delta_Q$ is the Laplacian attached to the dual quadratic
form.  In the degrees used below, $\mathcal H_d$ is the top harmonic
summand, while powers of $Q$ give radial summands.

\begin{lemma}[The edge-selected alternating cycle]
\label{lem:shared-radial-cycle}
For \(q=7,11\), let \(\Omega_+\sqcup\Omega_-\) be the two
one-factorization sheets.  Every endpoint edge \(e\) lies in a unique
pair \(M_+(e)\in\Omega_+\), \(M_-(e)\in\Omega_-\).  The matchings share
only \(e\), and after deleting it their union is one alternating cycle
on the other \(q-1\) endpoints.  The corresponding conic quotient has
nonzero deepest radial projection.
\end{lemma}

\begin{proof}
Move \(e\) to \(\{0,\infty\}\), so its secant is \(Y\), and put
\[
 n=\frac{q-1}{2},\qquad Q_0=(\F_q^\times)^2.
\]
The complementary matchings have the common torus normal form
\[
 \{\{u,cu\}:u\in Q_0\},\qquad
 \{\{u,c^{-1}u\}:u\in Q_0\},                         \tag{3.3a}
\]
where \(c^2=1/4\).  Here is the promised derivation.  The setwise edge
stabilizer in \(G^+\) is
\[
 D_e=\{x\mapsto sx:s\in Q_0\}
 \sqcup\{x\mapsto rs/x:s\in Q_0\},
\]
where \(r\) is any nonsquare.  It has order \(q-1\) and acts regularly on
\(\F_q^\times\): neither a nonidentity multiplication nor an inversion of
the displayed form has a fixed point.  Since each sheet is a
one-factorization, its unique matching through \(e\) is \(D_e\)-invariant.

One labeled calculation identifies the relevant involutions.  For \(q=7\),
\(x\mapsto x/(2-x)\) sends the edge \(\{0,2\}\) of \(M_{B_3}\) to \(e\) and
the complement to
\[
 \{\{1,5\},\{2,3\},\{4,6\}\}
 =\{\{u,5u\}:u\in Q_0\}.
\]
For \(q=11\), \(x\mapsto x/(1-x)\) sends \(\{0,1\}\) of \(M_{H_3}\) to
\(e\) and the complement to
\[
 \{\{1,2\},\{3,6\},\{4,8\},\{5,10\},\{7,9\}\}
 =\{\{u,2u\}:u\in Q_0\}.
\]
Choose \(c=3\) for \(q=7\) and \(c=6\) for \(q=11\).  In each case \(c\)
is the nonsquare root of \(c^2=1/4\), and the displayed complement is the
\(c^{-1}\)-matching.  Multiplication by any nonsquare fixes \(e\), exchanges
the two sheets, and rewrites that matching as the \(c\)-matching.  This
proves (3.3a), with \(c\) specifically nonsquare.  Multiplication by
\(c^{-2}=4\) has order \(n\) on \(Q_0\) for \(q=7,11\), proving the
alternating-cycle assertion.

Write
\[
 A_c(X,Y,Z)=
 \prod_{u\in Q_0}\bigl(cu^2X-(1+c)uY+Z\bigr).
\]
The two matching products share the factor \(Y\), and equality on the
conic gives
\[
 A_{c^{-1}}-A_c=(XZ-Y^2)\Psi_c.                         \tag{3.3b}
\]
If
\[
 h_j=\frac{n}{n-j}\binom{n-j}{j},
\]
the resultant with \(u^n-1\) gives one Dickson recurrence:
\[
 A_{c^{-1}}-A_c=
 2\sum_{j=0}^{(n-1)/2}
 (-1)^j h_jc^j(1+c)^{n-2j}(XZ)^jY^{n-2j}.              \tag{3.3c}
\]
For \(q=7\), division in (3.3b) gives
\(\Psi_c=5Y\).  For \(q=11\), it gives
\[
 \Psi_c=2Y^3+XZY,\qquad
 \Delta_Q\Psi_c=3Y.
\]
The full quotient is \(Y\Psi_c\), up to the fixed nonzero normalization.
The common secant \(Y\) has nonzero dual conic norm, so the displayed
linear deepest traces force a nonzero radial Fischer component in both
degrees.
\end{proof}

\begin{remark}[The Paley carrier]
\label{rem:paley-carrier}
Label either sheet by its regular translation subgroup.  The cross-sheet
incidence matrix is the circulant support of
\(\{0\}\cup(\F_q^\times)^2\), the complement of the Paley difference
set.  Its bordered sign matrix is the Paley Hadamard matrix of order
\(q+1\).  This explains the multiplier \(4\) and the sheet orientation,
but it is not the Gorenstein pairing used below.
\end{remark}

\begin{theorem}[The \(3,6,10\) quotient ranks]
\label{thm:rank-three-quotients}
For the matching configurations in (3.1), let
\[
 W_T=\operatorname{span}_{\F_q}
 \{\Phi_{M_T}(M):M\in\Omega_T\}.
\]
Then
\[
\begin{array}{c@{\quad}c@{\quad}c@{\quad}c}
T&\deg\Phi& W_T&\dim W_T\\ \hline
A_3&1&R_1=\mathcal H_1&3,\\
B_3&2&R_2=\mathcal H_2\oplus\F_7Q&6,\\
H_3&4&\mathcal H_4\oplus\F_{11}Q^2&10.
\end{array}                                               \tag{3.4}
\]
If \(h_T=4,6,10\) is the Coxeter number and
\(d=h_T/2-1\), the three rows have the uniform form
\[
 W_T=
 \begin{cases}
 \mathcal H_d,&d\ \text{odd},\\
 \mathcal H_d\oplus\F_qQ^{d/2},&d\ \text{even},
 \end{cases}
 \qquad
 \dim W_T=h_T-1+\mathbf 1_{\{d\ {\rm even}\}}.
\]
In particular, the first two configurations fill the complete quotient
form space.  The \(H_3\) configuration spans a canonical ten-dimensional
subspace of the fifteen-dimensional space \(R_4\).
\end{theorem}

\noindent\emph{Proof mode.}
The quotient construction and all three spanning assertions are proved
symbolically.  Irreducibility forces the \(A_3\) image and the top \(B_3\)
summand.  For \(H_3\), a Sylow-normalizer cohomology argument excludes the
middle Fischer layer and forces the top harmonic summand.  The common
edge-selected alternating cycle supplies both radial lines.  A separate
implementation reconstructs the three full quotient matrices and every
edge-selected cycle from the data in (3.1).

\begin{proof}
A matching in (3.1) has
$m=3,4,6$ edges, respectively, so
Proposition~\ref{prop:matching-secant-quotient} places its quotient in
$R_{m-2}$, of degrees $1,2,4$.  Direct differentiation gives
\[
 \dim\mathcal H_1=3,\qquad
 \dim\mathcal H_2=5,\qquad
 \dim\mathcal H_4=9.
\]
The radial vector \(Q\in R_2\) is not harmonic in characteristic \(7\).

We first prove the \(A_3\) and \(B_3\) rows without quotient matrices.
For \(q=5,7\), put
\[
 G=\operatorname{PGL}_2(q),\qquad
 G^+=\operatorname{PSL}_2(q),\qquad
 \chi(g)=\det(g)^{(q-1)/2}.
\]
The Dickson form \(s^qt-st^q\), which is the common conic restriction
of the matching products, has contragredient weight \(\det^{-1}\).
The Clebsch identification and Fischer decomposition therefore give
\[
\begin{array}{c@{\qquad}c}
q=5&\widetilde R_1=R_1\otimes\det^{-1}=M_2,\\
q=7&\widetilde R_2=R_2\otimes\det^{-1}=M_4\oplus M_0,
\end{array}                                                \tag{3.4a}
\]
where, uniformly,
\[
 M_r=\operatorname{Sym}^r(V^\vee)\otimes
       \det^{r/2}\otimes\chi,\qquad M_0=\chi.
\]
The scalar action is trivial, so these are \(G\)-modules.  On \(G^+\),
\(M_2\) and \(M_4\) are the absolutely irreducible symmetric powers of
highest weights \(2<5\) and \(4<7\), respectively
\cite[Fact~3.1]{BraunNebe2025}.

For either type let
\[
 c(g)=\Phi_{M_T}(gM_T).
\]
It is an affine cocycle, \(c(gh)=g c(h)+c(g)\), and its value span is
\(G\)-stable because
\[
 b\,c(g)=c(bg)-c(b).                                      \tag{3.4b}
\]
Distinct matchings have distinct products: unique factorization recovers
the normalized secant factors.  Hence every nonbase orbit point gives a
nonzero cocycle value.  In type \(A_3\), irreducibility of \(M_2\) now
forces the value span to be all of \(R_1\).

In type \(B_3\), project the cocycle to \(M_0=\chi\).  Its restriction
to \(G^+\) is a homomorphism from the simple group
\(\operatorname{PSL}_2(7)\) to the additive group of \(\F_7\), and is
therefore zero.  A nonbase point on the \(G^+\)-sheet consequently gives
a nonzero vector in \(M_4\); irreducibility and (3.4b) force the whole
five-space \(M_4=\mathcal H_2\).

Lemma~\ref{lem:shared-radial-cycle} supplies an outer-sheet value with
nonzero radial projection.  Since \(\Delta_Q\) kills
\(M_4=\mathcal H_2\), its value has nonzero
\(M_0=\F_7Q\)-projection.  The span already contains \(M_4\), so it
contains \(M_0\) as well.  Hence \(W_{B_3}=R_2\), proving the first two
rows of (3.4).

We prove the \(H_3\) row equivariantly.  Put
\(G=\operatorname{PGL}_2(11)\),
\(G^+=\operatorname{PSL}_2(11)\), and let
\(\chi(g)=\det(g)^5\) be the square-class character.  On the conic, the
common restriction of the matching products is the Dickson form
\(s^{11}t-st^{11}\), of contragredient weight \(\det^{-1}\).
Let \(V=\F_{11}^2\), with binary forms carrying the contragredient
\(\operatorname{GL}_2(11)\)-action.  The equivariant Fischer decomposition
obtained from the Clebsch identification
\(\operatorname{Sym}^2V\simeq\F_{11}^3\), together with the Dickson
weight in the quotient transformation (2.2), gives the linearized module
\[
 \widetilde R_4=R_4\otimes\det^{-1}
  =M_8\oplus M_4\oplus M_0.                       \tag{3.5}
\]
Explicitly, \(g\in\operatorname{GL}_2(11)\) acts on the first expression
by
\[
 g\cdot f=\det(g)^{-1}\bigl(f\circ\rho(g)^{-1}\bigr).
\]
The three summands are
\[
\begin{split}
 M_8&=\mathcal H_4\otimes\det^{-1}
      \simeq\operatorname{Sym}^8(V^\vee)\otimes\det^{-1},\\
 M_4&=Q\mathcal H_2\otimes\det^{-1}
      \simeq\operatorname{Sym}^4(V^\vee)\otimes\det^{-3},\\
 M_0&=\F_{11}Q^2\otimes\det^{-1}
      \simeq\det^{-5}=\chi .
\end{split}                                         \tag{3.6}
\]
Equivalently, all three summands have the uniform form
\[
 M_r=\operatorname{Sym}^r(V^\vee)\otimes
       \det^{r/2}\otimes\chi,\qquad r=8,4,0.
\]
A scalar \(\lambda I\) acts here as
\(\lambda^{-r}\lambda^r(\lambda^2)^5=1\), so these modules and their
direct sum descend from \(\operatorname{GL}_2(11)\) to \(G\).  This is
the \(G\)-action used below.  The determinant twist changes only the
linearization, not the underlying Fischer subspaces
\(\mathcal H_4,Q\mathcal H_2,\F_{11}Q^2\) of \(R_4\).

We need two elementary cohomology facts.  Let
\(U=\langle u\rangle\) be the upper-unipotent Sylow \(11\)-subgroup
and let \(a\) generate the diagonal torus of order \(10\).
Since \(r<10\), the cyclic norm
\(1+u+\cdots+u^{10}=(u-1)^{10}\) vanishes on \(M_r\), while
\((u-1)M_r\) has codimension one.  Thus \(H^1(U,M_r)\) is represented by
the class of \(x^r\).  Diagonal conjugation acts on that class by
\[
 a^{\,4-r/2}.                                      \tag{3.7}
\]
Restriction from \(G\) to \(U\) is injective because
\([G:U]=120\) is invertible in \(\F_{11}\), and its image is fixed by the
Sylow normalizer.  Equation (3.7) therefore gives
\[
 H^1(G,M_4)=H^1(G,M_0)=0,\qquad
 \dim H^1(G,M_8)\le1.                              \tag{3.8}
\]
The stabilizer \(H=A_5\) of \(M_{H_3}\) lies in \(G^+\).
The restrictions of \(M_8,M_4,M_0\) to \(H\) are respectively
\(4\oplus5,5,1\), so
\[
 M_8^H=M_4^H=0,\qquad \dim M_0^H=1.                \tag{3.9}
\]
Equivalently, average the characters
\((9,1,0,-1,-1)\), \((5,1,-1,0,0)\), and
\((1,1,1,1,1)\) over the \(A_5\) class sizes
\((1,15,20,12,12)\).

Let \(c(g)=\Phi_{M_{H_3}}(gM_{H_3})\).  The common-restriction identity
gives
\[
 c(gh)=g c(h)+c(g),
\]
and \(c\) vanishes on \(H\).  Project \(c\) through (3.5).  By (3.8), its
\(M_4\)-component is a coboundary; its translating vector lies in
\(M_4^H=0\), so that component vanishes.  This proves, without quotient
coordinates,
\[
 \Delta_Q\Phi_{M_{H_3}}(M)\in\F_{11}Q
 \quad\text{for every }M\in\Omega_{H_3}.            \tag{3.10}
\]
Likewise the \(M_0\)-component is a coboundary and hence vanishes on
\(G^+\).  Choose any other matching in the eleven-element \(G^+\)-orbit
of \(M_{H_3}\).  Its quotient is nonzero: equality of the two secant
products would identify their irreducible line factors and hence the
matchings.  Its \(M_4\)- and \(M_0\)-components vanish, so it gives a
nonzero vector in \(M_8\).  The degree-eight restricted
\(\operatorname{SL}_2(\F_{11})\)-module is irreducible
\cite[Fact~3.1]{BraunNebe2025}.  The span of the cocycle values is
\(G\)-stable, since
\[
 b\,c(g)=c(bg)-c(b)\qquad(b,g\in G).
\]
Therefore the quotient span contains all of \(M_8\).

Lemma~\ref{lem:shared-radial-cycle} supplies an outer-sheet value with
nonzero deepest radial projection.  Since \(\Delta_Q^2\) kills
\(M_8\oplus M_4\), this supplies the \(M_0\)-projection.  The span
already contains \(M_8\), so it also contains \(M_0\).  Hence
\(W_{H_3}=\mathcal H_4\oplus\F_{11}Q^2\), of dimension \(10\).
Since \(\dim R_4=15\), this equality is proper.
\end{proof}

\begin{remark}[Historical radial coordinates]
\label{rem:historical-radial-scalars}
Direct coordinates give \(\Delta_Q\Phi=4\) for the displayed \(B_3\)
pair and \(\Delta_Q^2\Phi=10\) for the displayed \(H_3\) pair.  These
values independently check the two specializations of
Lemma~\ref{lem:shared-radial-cycle}; they are not premises of its
Dickson-recurrence proof.
\end{remark}

\begin{corollary}[The base-choice class and its homogeneous extension]
\label{cor:h3-affine-origin}
Let \(c_8:G\to M_8\) be the top-harmonic projection of the base-choice
cocycle \(c\).  Then
\[
 H^1(G,M_8)=\F_{11}[c_8].                              \tag{3.11a}
\]
In particular, the affine \(H_3\) quotient configuration has no
\(G\)-equivariant origin: no translation makes its affine action linear.
Equivalently, the homogeneous module
\[
 E_{c_8}=M_8\oplus\F_{11},\qquad
 g\cdot(v,t)=(gv+t\,c_8(g),t),                         \tag{3.11b}
\]
is nonsplit.  Up to equivalence and nonzero rescaling of the quotient
coordinate, it is the unique nonsplit extension of the trivial
\(G\)-module by \(M_8\).
\end{corollary}

\begin{proof}
Equation (3.8) gives \(\dim H^1(G,M_8)\leq1\).  On the base sheet the
\(M_0\)-component vanishes, and the \(M_4\)-component vanishes by (3.9);
the nonzero same-sheet quotient used in the proof above is therefore a
value of \(c_8\).  If \(c_8(g)=gv-v\), its vanishing on
\(H=\operatorname{Stab}_G(M_{H_3})\) would put \(v\) in \(M_8^H=0\),
forcing \(c_8=0\), a contradiction.  Thus \([c_8]\ne0\), proving
(3.11a).

Changing the reference quotient translates the affine configuration
and changes \(c_8\) by a coboundary.  Hence a \(G\)-equivariant origin
would kill \([c_8]\), which is impossible.  Finally the standard
identification
\(\operatorname{Ext}^1_{\F_{11}G}(\F_{11},M_8)\simeq H^1(G,M_8)\)
identifies (3.11b) with \([c_8]\); the one-dimensionality just proved
gives the last assertion.
\end{proof}

\begin{corollary}[The exact \(H_3\) loss]
\label{cor:h3-middle-layer}
Over \(\F_{11}\), the quartic form space has the Fischer decomposition
\[
 R_4=\mathcal H_4\oplus Q\mathcal H_2\oplus\F_{11}Q^2.
                                                               \tag{3.12}
\]
Consequently,
\[
 R_4=W_{H_3}\oplus Q\mathcal H_2,\qquad
 R_4/W_{H_3}\simeq Q\mathcal H_2.                              \tag{3.13}
\]
Equivalently,
\[
 W_{H_3}
 =\{f\in R_4:\Delta_Qf\in\F_{11}Q\}.                           \tag{3.14}
\]
Thus the \(H_3\) quotient configuration retains the top harmonic
nine-space and the deepest radial line, and omits precisely the
five-dimensional middle radial layer.
\end{corollary}

\begin{proof}
For a homogeneous form \(f\) of degree \(r\), direct differentiation
gives
\[
 \Delta_Q(Qf)=(4r+6)f+Q\Delta_Qf.                              \tag{3.15}
\]
Hence \(\Delta_Q(Qh)=3h\) for \(h\in\mathcal H_2\), while
\(\Delta_Q(Q^2)=9Q\) over \(\F_{11}\).  Applying
\(\Delta_Q^2\) and then \(\Delta_Q\) shows that the three summands in
(3.12) meet trivially.  Their dimensions are \(9,5,1\), so they exhaust
the fifteen-dimensional space \(R_4\).  Theorem
\ref{thm:rank-three-quotients} then gives (3.13).  If
\(f=h_4+Qh_2+cQ^2\) is its decomposition, then
\[
 \Delta_Qf=3h_2+9cQ.
\]
The quadratic decomposition \(R_2=\mathcal H_2\oplus\F_{11}Q\)
therefore gives (3.14).
\end{proof}

The rank proof has one general part and one exceptional input.  The conic
ideal gives the common restriction and quotient, while irreducibility
forces the \(A_3\) image and the top \(B_3\) summand.  For \(H_3\),
(3.7)--(3.10) use the specific
\(A_5<\operatorname{PGL}_2(11)\) module data to force the top summand and
exclude the middle one.  Lemma~\ref{lem:shared-radial-cycle} supplies both
radial lines from the same recurrence.  Appendix A records the detector
calculations and their finite seams.

Corollary~\ref{cor:h3-middle-layer} identifies the linear information
lost by the \(H_3\) orbit, but linear span does not recover the two
factorization sheets inside the finite configuration.  Their recovery
begins one categorical level higher, with products of affine evaluation
functions.

\section{Balanced sheets and cubic orientation}
\label{sec:balanced}

For \(T=B_3,H_3\), put \(q=7,11\), respectively.  The subgroup
\[
 G^+=\operatorname{PSL}_2(q)
 \subset G=\operatorname{PGL}_2(q)
\]
splits \(\Omega_T\) into two orbits
\(\Omega_+\sqcup\Omega_-\), each of size \(q\); an element of the outer
coset exchanges them.  These are the two one-factorization sheets.  The
\(A_3\) orbit does not split: its \(S_4\) stabilizer cannot lie in
\(\operatorname{PSL}_2(5)\simeq A_5\), already by order.

Write \(x_M=\Phi_{M_T}(M)\in W_T\).  Let \(L\subseteq\F_q^{\Omega_T}\)
be the evaluation space of affine-linear functions on the finite
configuration \(\{x_M\}\).  Since the points affinely span \(W_T\),
\[
 \dim L=1+\dim W_T=q.
\]
For subspaces of functions, write
\[
 L^{\circ2}=\operatorname{span}\{fg:f,g\in L\},
\]
where multiplication is pointwise.
The common second-moment form in the proposition below is
\[
 B(\lambda,\mu)=\sum_{M\in\Omega_+}\lambda(x_M)\mu(x_M)
               =\sum_{M\in\Omega_-}\lambda(x_M)\mu(x_M)
\]
on \(W_T^\vee\).  A \emph{field-valued strength-two trade} is a
weight \(\tau\in k^\Omega\) satisfying, for all \(f,g\in L\),
\[
 \sum_{M\in\Omega}\tau(M)f(M)g(M)=0.
\]

\paragraph{The \(B_3\) model.}
Here the quotient configuration consists of fourteen points spanning
the six-space \(W_{B_3}\), split into two seven-point sheets.  Its
affine evaluation space \(L\) has dimension seven.  Restriction to
either sheet has rank six, while the common second-moment form has rank
five and a radical functional that takes distinct constant values on
the two sheets.  It follows that \(L^{\circ2}\) is the
thirteen-dimensional hyperplane of functions with equal sheet sums;
its one-dimensional annihilator is the sheet sign.  Thus the quadratic
algebra recovers the \(7+7\) partition without a search through
partitions.  The first nonzero signed moment is cubic, so
\(L^{\circ3}=\F_7^{14}\), and the homogenized points have
\(h\)-vector \((1,6,6,1)\).  The next proposition isolates the
quadratic step in a form that applies to both balanced configurations.

\begin{proposition}[Radical--Hadamard recovery]
\label{prop:radical-hadamard}
Let \(k\) be a field and
\(\Omega=\Omega_+\sqcup\Omega_-\), with
\(|\Omega_+|=|\Omega_-|=q\) and \(q\cdot1_k=0\).  Suppose an affine evaluation space
\(L\subseteq k^\Omega\) has the following properties:
\begin{enumerate}[label=\textup{(\roman*)},leftmargin=*]
\item each sheet has zero first moment, and the two sheets have equal
second moments;
\item \(\dim L=q\), and restriction from \(L\) onto the zero-sum
hyperplane of either sheet has rank \(q-1\);
\item the second-moment form on linear functionals has rank \(q-2\)
and a one-dimensional radical;
\item a nonzero radical functional is constant on each sheet, with
different constants on the two sheets.
\end{enumerate}
Then
\[
 L^{\circ2}
 =\left\{(u_+,u_-):
   \sum_{\Omega_+}u_+=\sum_{\Omega_-}u_-\right\},              \tag{4.1}
\]
and
\[
 (L^{\circ2})^\perp
 =k(\mathbf1_{\Omega_+}-\mathbf1_{\Omega_-}).                 \tag{4.2}
\]
Consequently every field-valued strength-two trade is a scalar multiple
of the sheet sign.
\end{proposition}

\begin{proof}
Let \(e_+\) and \(e_-\) be the two sheet indicators.  The evaluation
vector of a nonzero radical functional is \(a_+e_++a_-e_-\), with
\(a_+\ne a_-\).
Together with \(1=e_++e_-\in L\), it therefore yields
\(e_+,e_-\in L\).

Let \(H_\pm\) be the zero-sum hyperplane on \(\Omega_\pm\).
The first-moment hypothesis and \(q\cdot1_k=0\) put each restriction of \(L\)
inside \(H_\pm\); the rank hypothesis makes both restrictions
surjective.  Multiplication by \(e_+\) and \(e_-\) now gives
\[
 (H_+\oplus0)+(0\oplus H_-)\subseteq L^{\circ2}.              \tag{4.3}
\]
This subspace has dimension \(2q-2\).

For \(f,g\in L\), the difference between the two sheet sums of \(fg\)
expands into constant, first-moment, and second-moment terms.  The
hypotheses make all three differences zero, so \(L^{\circ2}\) lies in
the hyperplane on the right of (4.1).  Finally choose restrictions
\(a,b\in H_+\) with nonzero standard pairing and lift them to
\(f,g\in L\).  Equality of the second moments gives the same nonzero
pairing on \(\Omega_-\).  Thus \(fg\) has equal nonzero sheet sums and
does not lie in (4.3).  We have found \(2q-1\) independent directions
in the \(2q-1\)-dimensional hyperplane, proving (4.1).  Its annihilator
under the standard coordinate pairing is the line in (4.2).
\end{proof}

\begin{proposition}[The modular-permutation mechanism]
\label{prop:modular-sheet-mechanism}
For \(T=B_3,H_3\), the affine quotient configuration satisfies all four
hypotheses of Proposition~\ref{prop:radical-hadamard}.  More precisely,
write
\[
 W_T=U_T\oplus\F_q r,\qquad
 U_{B_3}=\mathcal H_2,\quad U_{H_3}=\mathcal H_4.
\]
On either \(q\)-point sheet, affine evaluation has image the zero-sum
hyperplane and rank \(q-1\).  The common second-moment form has rank
\(q-2\), with radical the radial functional \(r^\vee\); that functional
is constant on each sheet and its two constants are distinct.
\end{proposition}

\begin{proof}
As a \(G^+\)-module, \(U_T\) is the irreducible module
\(M_{q-3}\) of dimension \(q-2\).  Project the affine quotient points to
\(U_T\), writing the resulting sheet as \(\{y_\omega\}\).
The cocycle argument in Section~\ref{sec:rank-three} says that the
differences \(y_\omega-y_{\omega'}\) span \(U_T\).
Moreover \(\sum_\omega y_\omega\) is \(G^+\)-fixed: translating the
sheet adds \(q\) times the cocycle value, which vanishes in
characteristic \(q\).  Since \(U_T^{G^+}=0\),
\[
 \sum_\omega y_\omega=0.
\]

Let \(P=\F_q^{\Omega_+}\), and put
\[
 P_0=\left\{(a_\omega):\sum_\omega a_\omega=0\right\},
 \qquad S=\F_q\mathbf1.
\]
Here \(S\subset P_0\), the defining-characteristic feature that drives
the argument.  The map
\[
 P_0\longrightarrow U_T,\qquad
 (a_\omega)\longmapsto\sum_\omega a_\omega y_\omega
\]
is onto because the differences span, and the preceding moment identity
kills \(S\).  Both \(P_0/S\) and \(U_T\) have dimension \(q-2\), so it
induces an isomorphism
\[
 P_0/S\simeq U_T.
\]
The standard dot product on \(P_0\) has radical exactly \(S\).
Dualizing therefore shows that the evaluations of \(U_T^\vee\) form a
nondegenerate \((q-2)\)-space complementary to \(S\) in \(P_0\).
Adding constants gives all of \(P_0\), proving the restriction rank and
the top second-moment rank.  The same argument applies to the other
sheet.

The two top second moments are equal.  On
\[
 M_{q-3}=\operatorname{Sym}^{q-3}(V^\vee)
          \otimes\det^{(q-3)/2}\otimes\chi
\]
the standard apolar pairing is \(G\)-invariant: its determinant weight
\(\det^{-(q-3)}\) is cancelled by the square of the displayed twist,
and \(\chi^2=1\).  Each sheet moment is a nonzero
\(G^+\)-invariant form, hence a scalar multiple of this pairing, while
an outer element transports one sheet to the other by an apolar
isometry.  Their scalars, and thus the two forms, are equal.

Finally, projection of the affine cocycle to the radial character
\(\chi\) restricts on the simple group \(G^+\) to an additive
homomorphism, hence vanishes.  The radial coordinate is consequently
zero on the base sheet and constant on the other.  The secant-norm
calculation in Lemma~\ref{lem:shared-radial-cycle} shows that the outer
constant is nonzero in both types.  Because
\(q\cdot1_{\F_q}=0\), radial--radial and radial--top second moments vanish.  Thus
\(r^\vee\) is exactly the one-dimensional radical, and its sheet
constants are distinct.
\end{proof}

\Needspace{15\baselineskip}
Let \(T=B_3,H_3\), with \(q=7,11\), respectively, and fix a reference
matching \(M_T\) with stabilizer \(K=S_4,A_5\).  Put
\(m=(q+1)/2\) and \(d=m-2\).  Write \(\mathcal A_T\) for the affine space
of degree-\(m\) forms whose restriction to the conic is the fixed binary form
vanishing on all \(q+1\) marked rational points.  We call its members that
split completely into linear factors over \(\overline{\F}_q\) the
\emph{geometric Chow locus}.  All intersections below are set-theoretic.
Symmetry leaves one affine parameter in \(\mathcal A_T\); geometric
splitting selects one rational point.

\begin{theorem}[The fixed line and its Chow point]
\label{thm:fixed-line-chow-rigidity}
With the preceding notation and the determinant normalization of
Section~\ref{sec:matching-products}, the following hold.

\begin{enumerate}[label=\textup{(\roman*)},leftmargin=*]
\item The \(K\)-fixed locus is an affine line:
\[
 \mathcal A_T^K
   =P_{M_T}+Q(R_d)^K
   =P_{M_T}+\F_q Q^{d/2+1}.                         \tag{4.4a}
\]
Its \(q\) rational points lie in distinct \(G\)-orbits, each with
stabilizer \(K\).

\item The rational fixed line meets the geometric Chow locus in the single
point \(P_{M_T}\).  Equivalently, the matching point is its unique
geometrically split point.

\item The orbit of each noncoalescent point has \(2q\) points, and its
quadratic-trade space is the sheet-sign line.  Consequently \(q-2\) of
these orbits satisfy the same one-dimensional two-valued strength-two trade
condition as the matching orbit but do not come from matchings.
\end{enumerate}
Thus the quadratic-trade condition alone does not classify the
\(B_3/H_3\) configurations; the matching carrier in
Theorem~\ref{thm:balanced-orbit-completeness} is essential.  Within that
carrier, the classification remains exact.
\end{theorem}

\begin{proof}
The translation space of \(\mathcal A_T\) is \(QR_d\), and \(P_{M_T}\)
is \(K\)-fixed.  Since \(q\nmid|K|\), invariant dimensions may be computed
by averaging characters.  For \(B_3\), the invariant subspace of the quadratic
form module under the irreducible orthogonal \(S_4\)-action is
\(\F_7Q\).  For \(H_3\), the character of the quartic form module on the
classes of orders \(1,2,3,5,5\) has values \(15,3,0,0,0\); hence
\[
 \dim R_4^{A_5}=\frac{15+15\cdot3}{60}=1,
\]
and \(R_4^{A_5}=\F_{11}Q^2\).  This proves (4.4a).

The proof of Theorem~\ref{thm:balanced-orbit-completeness} shows that
\(N_G(K)=K\).  A point on (4.4a) cannot be fixed by \(G^+\): its
difference from \(P_{M_T}\) is radial and hence \(G^+\)-fixed, which
would make \(P_{M_T}\) itself \(G^+\)-fixed.  Maximality of \(K\) in
\(G^+\), together with \(N_G(K)=K\), therefore makes its stabilizer
exactly \(K\).  If two points of the line were \(G\)-conjugate, a
conjugating element would normalize \(K\), hence lie in \(K\), and the
points would be equal.  This proves the last assertion of (i).

Suppose a \(K\)-fixed point of \(\mathcal A_T\) is a product of linear
forms.  Each factor cuts a degree-two divisor on the conic, while their
product cuts the squarefree divisor of all \(q+1\) marked points.  Thus
the factors pair those points into marked secants.  This remains true when
the factorization is taken over \(\overline{\F}_q\): every factor divisor is
a two-element subset of the displayed rational divisor, so its joining line
is rational.  Unique
factorization makes that matching \(K\)-invariant.  For \(S_4\) on eight
points, a point stabilizer is
\(C_3\) and the unique possible two-point block stabilizer is
\(N_{S_4}(C_3)=S_3\).  For \(A_5\) on twelve points, the corresponding
groups are \(C_5\) and \(N_{A_5}(C_5)=D_{10}\).  Thus in either case the
\(K\)-invariant matching is unique, proving (ii).

It remains to track the quadratic relation, not the individual orbits.
Write the fixed line in quotient coordinates as \(x_t=x_0+tr\), where
\(r=Q^{d/2}\) spans the radial character.  On \(G^+\) this character is
trivial, while the outer coset negates it.  Relative to the moving base
point \(x_t\), the top harmonic configurations on both sheets are
therefore unchanged; one radial sheet constant remains \(0\), and the
other changes from the nonzero constant \(c\) of
Proposition~\ref{prop:modular-sheet-mechanism} to \(c-2t\).
Because \(q\cdot1_{\F_q}=0\) and each top sheet has zero sum, this
translation changes neither the first moments nor the two equal second
moments.

There is exactly one coalescence parameter \(t=c/2\).  For every other
\(t\), all four hypotheses of
Proposition~\ref{prop:radical-hadamard} are unchanged, so the
quadratic-trade space is exactly the sheet-sign line.  The matching
point is \(t=0\), and \(c\ne0\), so it is not the coalescence point.
Among the \(q-1\) noncoalescent orbits, exactly one is the matching
orbit by (ii); the remaining \(q-2\) prove (iii).
\end{proof}

\begin{remark}[The coalescence point]
\label{rem:fixed-line-coalescence}
At \(t=c/2\), the two radial constants agree and hypothesis~(iv) of
Proposition~\ref{prop:radical-hadamard} fails.  The square rank at this
single point is not needed for the classification or for
Theorem~\ref{thm:fixed-line-chow-rigidity}.
\end{remark}

For reference, exact reconstruction of the two configurations gives
\[
\begin{array}{c@{\quad}c@{\quad}c@{\quad}c@{\quad}c}
T&q&\text{sheet restriction ranks}
 &\text{second-moment rank}&\dim L^{\circ2}\\ \hline
B_3&7&6,6&5&13,\\
H_3&11&10,10&9&21.
\end{array}                                                   \tag{4.4}
\]
In each row the radical is one-dimensional and takes two distinct
sheet values.  Proposition~\ref{prop:modular-sheet-mechanism}, rather
than this table, proves these facts.  The exact certificate and its
independent reconstruction are cross-checks only; no enumeration of
equal-size subsets or row reduction is used in the proof.

\begin{lemma}[A full-support quadratic relation forces the cube]
\label{lem:hyperplane-square}
Let \(\Omega\) be finite and let \(L\subseteq k^\Omega\) be a unital
linear space with \(\dim L>1\).  Suppose
\[
 L^{\circ2}
 =\left\{u:\sum_{\omega\in\Omega}\epsilon_\omega u_\omega=0\right\}
\]
for some \(\epsilon\in(k^\times)^\Omega\).  Then
\[
 L^{\circ3}=k^\Omega.
\]
\end{lemma}

\begin{proof}
If a nonzero vector \(\eta\) annihilated \(L^{\circ3}\), then
\(L^{\circ2}\subseteq L^{\circ3}\) would put \(\eta\) in
\(k\epsilon\).  For every \(f\in L\), the vector
\(\eta\circ f\) also annihilates \(L^{\circ2}\), because
\[
 \langle\eta\circ f,h\rangle=\langle\eta,fh\rangle=0
 \qquad(h\in L^{\circ2}).
\]
Thus \(\eta\circ f\in k\epsilon=k\eta\).  Since \(\eta\) has full
support, every \(f\in L\) is constant, contrary to \(\dim L>1\).
The annihilator of \(L^{\circ3}\) is therefore zero.
\end{proof}

\begin{theorem}[Balanced sheets and cubic-first orientation]
\label{thm:balanced-cubic}
For \(T=B_3,H_3\), the unordered pair
\(\{\Omega_+,\Omega_-\}\) is intrinsic to the affine quotient
configuration:
\begin{enumerate}[label=\textup{(\roman*)},leftmargin=*]
\item up to interchange, \(\Omega_+\) and \(\Omega_-\) are the only
complementary halves with equal first and second tensor moments;
\item if \(\epsilon\) is \(+1\) on \(\Omega_+\) and \(-1\) on
\(\Omega_-\), then
\[
 \mu_j=\sum_{M\in\Omega_T}\epsilon(M)x_M^{\odot j}
       \in\operatorname{Sym}^j(W_T)
\]
satisfies
\[
 \mu_1=0,\qquad \mu_2=0,\qquad \mu_3\ne0;                    \tag{4.5}
\]
\item \(\mu_3\) is independent of the reference matching,
\(G^+\) fixes it, and the outer coset negates it.  Within \(G\),
\[
 \operatorname{Stab}_G(\mu_3)=G^+,\qquad
 \operatorname{Stab}_G(\F_q\mu_3)=G.                         \tag{4.6}
\]
\end{enumerate}
Thus degree three is the first signed tensor moment, or linear
functional of such a moment, that orients the recovered sheet pair.
This minimality statement concerns signed tensor moments, not every
possible nonlinear statistic of the configuration.
\end{theorem}

\noindent\emph{Proof mode.}
Proposition~\ref{prop:radical-hadamard} is a symbolic proof.
Proposition~\ref{prop:modular-sheet-mechanism} proves its hypotheses from
the defining-characteristic permutation module and the shared
alternating-cycle witness.  Lemma~\ref{lem:hyperplane-square} forces the
nonzero cubic without a tensor calculation or Gorenstein premise.  Signed
Gale duality records self-association, while maximal-isotropic quotient
duality gives the Gorenstein pairing.  Exact finite arithmetic
independently cross-checks the displayed ranks.

\begin{proof}
Equality of first and second tensor moments for a complementary pair is
equivalent, by polarization, to its sign weight annihilating
\(L^{\circ2}\).  Since the characteristics are \(7\) and \(11\),
polarization through degree two is valid.  Proposition
\ref{prop:radical-hadamard} says that the annihilator is exactly the
sheet-sign line.  A sign weight taking only the values \(\{\pm1\}\) is
therefore \(\epsilon\) or \(-\epsilon\), proving (i).

The same moment equalities give \(\mu_1=\mu_2=0\).
Proposition~\ref{prop:radical-hadamard} identifies
\(L^{\circ2}=\ker\langle\epsilon,-\rangle\), and \(\epsilon\) has full
support.  Lemma~\ref{lem:hyperplane-square} gives
\(L^{\circ3}=\F_q^{\Omega_T}\).  Hence \(\epsilon\) cannot annihilate all
cubic products.  Since the lower signed moments vanish and \(6\) is
invertible in \(\F_7,\F_{11}\), polarization gives \(\mu_3\ne0\).

For the geometric consequence, let \(A\) be the \(q\)-by-\(2q\) matrix
whose \(M\)-th column is
\(\widehat x_M=(1,x_M)\), and put
\(D=\operatorname{diag}(\epsilon(M))\).  Equal sheet sizes and the
vanishing signed moments through degree two give
\[
 A D A^{\mathsf T}=0.
\]
The affine spanning statement gives \(\operatorname{rank}A=q\), so the
rows of \(AD\) form the kernel of \(A\).  Thus \(AD\) is a Gale matrix,
and its columns differ from those of \(A\) only by the nonzero scalars
\(\epsilon(M)\).  The homogenized configuration is self-associated.

We now prove the Gorenstein statement directly.  On
\[
 V=\F_q^{\Omega_+}\oplus\F_q^{\Omega_-}
\]
use the nondegenerate signed coordinate form
\[
 B(u,v)=\sum_{\Omega_+}u_Mv_M-\sum_{\Omega_-}u_Mv_M.
                                                               \tag{4.6a}
\]
Proposition~\ref{prop:radical-hadamard} identifies
\[
 L^{\circ2}
 =\langle\mathbf1\rangle^{\perp_B}.                          \tag{4.6b}
\]
For \(u,v\in L\), the product \(u\circ v\) lies in the right side of
(4.6b), so \(B(u,v)=0\).  Thus \(L\) is isotropic.  Since
\(\dim L=q=\frac12\dim V\), it is maximal isotropic:
\[
 L=L^{\perp_B}.                                              \tag{4.6c}
\]

Let \(\overline R\) be the reduction of the homogeneous coordinate ring by
the homogenizing coordinate.  Its positive graded pieces are
\[
 \overline R_1=L/\langle\mathbf1\rangle,\qquad
 \overline R_2=L^{\circ2}/L,\qquad
 \overline R_3=V/L^{\circ2}.
\]
The form \(B\) descends to
\[
 \overline R_1\times\overline R_2\longrightarrow\F_q.        \tag{4.6d}
\]
Its left radical is
\((L^{\circ2})^\perp_B/\langle\mathbf1\rangle=0\), and its
right radical is \(L^\perp_B/L=0\).  Hence (4.6d) is perfect.
Lemma~\ref{lem:hyperplane-square} gives
\(L^{\circ3}=V\), so \(\overline R_3\) is one-dimensional and
multiplication realizes (4.6d) as the degree-\(1\)-by-degree-\(2\)
socle pairing.  Therefore \(\overline R\) is Artinian Gorenstein with
Hilbert function
\[
 (1,q-1,q-1,1).
\]
The homogenizing coordinate is regular because it is nonzero at every
point, so the projective coordinate ring is arithmetically Gorenstein.
The unique quadratic dependence is \(\epsilon\), which has full support;
hence every one-point deletion imposes independent conditions on
quadrics.

If the reference matching changes, every point translates
by one vector \(c\).  Expanding the signed third moment after
\(x_M\mapsto x_M+c\) leaves \(\mu_3\) unchanged: the other terms contain
\(\mu_2,\mu_1\), or
\(\mu_0=\sum_M\epsilon(M)=0\).

Every element of \(G^+\) preserves both sheets and hence fixes the signed
sum.  Every element of the outer coset exchanges the sheets and negates
it.  Since \(\mu_3\ne0\), these two transformation laws give (4.6).
Finally (4.5) proves the stated minimality inside the signed tensor-moment
family.
\end{proof}

\begin{corollary}[Graded evaluation algebra]
\label{cor:graded-evaluation}
Put \(L^{\circ0}=\langle1\rangle\), and let \(L^{\circ d}\) be the
span of the pointwise products of \(d\) elements of \(L\).
For \(T=B_3,H_3\), the pointwise powers of the affine evaluation
space satisfy
\[
 \dim L^{\circ d}=
 \begin{cases}
  1,&d=0,\\
  q,&d=1,\\
  2q-1,&d=2,\\
  2q,&d\ge3.
 \end{cases}
\]
In particular, \(L^{\circ3}=\F_q^{\Omega_T}\).  The homogenized
configuration
\[
 \widehat X_T=\{[1:x_M]:M\in\Omega_T\}\subseteq\mathbb P^{q-1}
\]
therefore has Hilbert function
\[
 1,\ q,\ 2q-1,\ 2q,\ 2q,\ldots
\]
and \(h\)-vector \((1,q-1,q-1,1)\).
\end{corollary}

\begin{proof}
Pair functions on \(\Omega_T\) with the sheet sign by
\(\langle\epsilon,f\rangle=\sum_M\epsilon(M)f(M)\).
Proposition~\ref{prop:radical-hadamard} gives
\(L^{\circ2}=\ker\langle\epsilon,-\rangle\), with \(\epsilon\) of full
support.  Lemma~\ref{lem:hyperplane-square} gives
\(L^{\circ3}=\F_q^{\Omega_T}\).  Multiplication by \(1\) gives the same
equality in every higher degree.  Homogeneous degree-\(d\) forms on
\(\widehat X_T\) restrict to \(L^{\circ d}\), which gives the Hilbert
function; its first difference is the stated \(h\)-vector.
Changing the reference matching translates every \(x_M\).  Such a
translation preserves \(L\) and all its pointwise powers and acts
projectively on \(\widehat X_T\), so the graded algebra and Hilbert
function are independent of the reference.
\end{proof}

\begin{corollary}[Self-association and cubic duality]
\label{cor:self-associated-gorenstein}
For \(T=B_3,H_3\), the configuration
\(\widehat X_T\subseteq\mathbb P^{q-1}\) is a reduced
self-associated arithmetically Gorenstein set of \(2q\) points.
Every one-point deletion imposes independent conditions on quadrics,
so \(\widehat X_T\) has the Cayley--Bacharach property in degree two.
Its socle residue vector is the sheet sign \(\epsilon\).  The
Artinian reduction by the homogenizing coordinate has Hilbert function
\((1,q-1,q-1,1)\), and its Macaulay inverse system is the cubic line
\(\F_q\mu_3\).
\end{corollary}

\begin{proof}
The proof of Theorem~\ref{thm:balanced-cubic} already establishes signed
Gale self-duality, the Cayley--Bacharach property in degree two, the
Gorenstein coordinate ring, and the displayed \(h\)-vector.  These
arguments commute with base change to an algebraic closure, so the
Gorenstein conclusion descends to \(\F_q\).

Let \(z\) be the homogenizing coordinate.  It is nonzero at every
point and is therefore regular on the coordinate ring.  The degree
\(d\) part of the reduction modulo \(z\) is
\(L^{\circ d}/L^{\circ(d-1)}\), so its Hilbert function is the
\(h\)-vector from Corollary~\ref{cor:graded-evaluation}.  The functional
\[
 f\longmapsto\sum_M\epsilon(M)f(x_M)
\]
annihilates degrees below three and is nonzero in degree three.
Thus \(\epsilon\) is the socle residue vector.  Degree-three
polarization identifies its Macaulay dual generator with \(\mu_3\).
\end{proof}

After base change to an algebraic closure, the signed coordinate form can
be put in the standard form used for self-dual codes.  In that form, the
implication above is the zero-defect case of the modern self-dual-code
criterion: the Gorenstein defect is measured by the
defect of the Schur square, or equivalently by decomposability
\cite[Theorem~3.11 and Corollaries~3.12--3.13]{RodriguezPajaresEtAl2025}.  Here
\(\dim L^{\circ2}=2q-1\) already places both configurations in the
zero-defect case.  The ideal, deletion, socle, inverse-system, and
minimal-resolution calculations are also checked directly in the
verification bundle.
Self-association does not select an affine origin.  Translating all
\(x_M\) acts projectively on the columns of \(A\) and preserves the
signed Gale identity, in agreement with the nontrivial base-choice
obstruction behind the quotient construction.

The cubic remembers an orientation that the linear quotient does not.
Indeed, the signed linear relation
\[
 \sum_{M\in\Omega_+}x_M=\sum_{M\in\Omega_-}x_M
\]
already shows that the sheet sign is killed by the linear
quotient map.  Quadratic products recover the unordered
partition through (4.2), and the cubic supplies the sign exchanged when
the two recovered levels are swapped.

\begin{corollary}[Secant-product tensor syzygies]
\label{cor:secant-product-syzygies}
For \(T=B_3,H_3\), let \(m_Q:R_{m-2}\to R_m\) denote multiplication
by \(Q\).  Regarding the displayed powers as symmetric tensors, not
polynomial powers, the canonically scaled secant products satisfy
\[
\begin{aligned}
\sum_M\epsilon(M)P_M&=0,\\
\sum_M\epsilon(M)P_M^{\odot2}&=0,\\
\sum_M\epsilon(M)P_M^{\odot3}
  &=\operatorname{Sym}^3(m_Q)(\mu_3)\ne0.
\end{aligned}                                               \tag{4.7}
\]
\end{corollary}

\begin{proof}
Substitute \(P_M=P_{M_T}+Qx_M\) and expand.  Each sheet has
\(q\cdot1_{\F_q}=0\), and the signed quotient moments vanish
through degree two.  All terms below cubic degree cancel, leaving
\(\operatorname{Sym}^3(m_Q)(\mu_3)\) in degree three.  Since
\(\F_q[X,Y,Z]\) is a domain, multiplication by \(Q\) is injective;
over a field, the symmetric powers of an injective linear map are
injective.
\end{proof}

\section*{Conclusion}

Restriction to a conic erases the pairing of its marked points.  Dividing
differences of secant products by the conic equation retains enough
information to recover exactly two balanced matching orbits:
\(B_3/\F_7\) and \(H_3/\F_{11}\).  The proof derives their
one-factorization sheets from the two-valued quadratic trade and excludes
all other odd prime powers without a matching census.  The same quotient
construction gives the ranks \(3,6,10\); for \(H_3\), it retains the top
harmonic part and deepest radial line and loses exactly
\(Q\mathcal H_2\).

The matching hypothesis is the precise boundary.  Each surviving
stabilizer fixes an affine line of conic products.  Its only completely
reducible point is the matching product, although \(q-2\) other orbits
retain the same two-valued trade.  The trade therefore remembers the
homogeneous sheet module but not its embedding as a secant factorization;
the Chow condition restores that information.

On the matching orbit, quadratic products recover the two sheets without
orienting them.  Their first signed tensor moment is cubic, and this cubic
orients the recovered pair.  It is also the Macaulay inverse system of the
self-associated Gorenstein homogenization.  The general Schur-square
criterion explains Gorensteinness; the configuration-specific result is
the chain from quadratic trade to matching orbit, from matching orbit to
unordered sheets, and from unordered sheets to the cubic orientation.
In these configurations the degree filtration therefore separates the
information carried by the construction: quadratic data recognize and
recover the unordered factorization, while the first odd degree both
restores orientation and supplies the Macaulay dual generator.

\appendix

\medskip
\noindent\textbf{Appendix guide.}
Appendix A is the load-bearing specialist layer for the uniform modular
exclusion: it proves the three detector interfaces stated in the main
argument.  Appendices B, C, E, and F record \(H_3\) refinements that use a
chosen \(A_4\)-structure; Appendix D compares the arithmetic marker data of
all three configurations.  No canonical identification of the depth plane
with the relative-cubic plane is asserted.
Appendix G gives the verification architecture for all statements.
The optional finite refinements in Appendices B, D, and F are explicitly
computer-assisted propositions: their concise tables are article-level
outputs, while the complete machine-readable inputs, deterministic
algorithms, checksums, and independent replays are the stable supplement
listed in Appendix G.  They are not presented as unaided human proofs.

\section{Detector calculations for uniform exclusion}
\label{app:modular-detectors}

The first-pass outputs of this appendix are
Lemmas~\ref{lem:targeted-linear-detectors},
\ref{lem:opposite-parity-quadratic}, and
\ref{lem:affine-class-contraction}.  The calculations below prove only
those detector channels; no exhaustive description of the finite-group
socle is asserted.

Fix \(k=\overline{\F}_q\), put \(E=L(1)\), and write
\[
 \mathbf G=\operatorname{SL}_2(k),\qquad
 \Gamma=\operatorname{SL}_2(q),\qquad
 d=\frac{q-3}{2},\qquad F=\operatorname{Sym}^dL(2),
\]
and set \(M=\operatorname{Sym}^2F\) and
\(Y_1=L(1)\otimes L(1)^{(e)}\).
Steinberg's tensor-product theorem classifies the simple
\(\Gamma\)-modules used here
\cite[Theorems~1.1 and~7.4]{Steinberg1963}.  Put
\(D=\nabla(d)=\operatorname{Sym}^dE\).  Modular Hermite reciprocity gives
the convention-explicit chain
\[
 \operatorname{Sym}^d(\operatorname{Sym}^2E)
 \simeq\operatorname{Sym}^d(\operatorname{Sym}_2E)
 \simeq\operatorname{Sym}_2(\operatorname{Sym}^dE)
 \simeq\operatorname{Sym}^2D.                         \tag{A.1a}
\]
Here \(\operatorname{Sym}_2\) is the lower symmetric square; the first and
last identifications are symmetrizations, valid because \(2<p\)
\cite[Corollary~1.5]{McDowellWildon2022}.  Thus
\(M\simeq\operatorname{Sym}^2(\operatorname{Sym}^2D)\); the targeted
arguments below do not require a socle formula for the inner symmetric
square.

\subsection{The linear detector channels}
\label{app:linear-detectors}

We first isolate the finite-root seam.  Let \(S=L(q-1)\) or \(L(q-7)\),
and let \(\phi:S\to F\) be a \(\Gamma\)-map.  For
\(u(t):(X,Y)\mapsto(X,Y+tX)\), put
\[
 D(t)=u_F(t)\phi-\phi u_S(t).
\]
The root matrices on \(F\) have degree at most \(q-3\); those on the two
sources have degree at most \(q-1\) and \(q-7\), respectively.  Hence
every entry of \(D(t)\) has degree at most \(q-1\).  It vanishes at all
\(q\) elements of \(\F_q\), so it is zero.  Thus \(\phi\) commutes with
the full algebraic positive root group.  The finite Weyl element supplies
the negative root group, and these two root groups generate \(\mathbf G\).
Consequently, in precisely these two channels,
\[
 \operatorname{Hom}_{\Gamma}(S,F)
 =\operatorname{Hom}_{\mathbf G}(S,F).                 \tag{A.1}
\]
This degree argument is deliberately not applied to \(M\), whose root
degree is \(2q-6\).

For \(S=\operatorname{St}=L(q-1)\), the top weight of \(F\) is \(q-3\).
Equation (A.1) therefore gives
\[
 \operatorname{Hom}_H(\operatorname{St},F)=0.          \tag{A.2}
\]

Now suppose \(q\equiv3\pmod4\) and \(e>1\).  Then
\(p\equiv3\pmod4\), \(e\) is odd, and every base-\(p\) digit of
\(q-7\) is even: from the least significant digit upward they are
\[
 (p-7,p-1,\ldots,p-1)\quad(p\ge7),\qquad
 (2,0,2,\ldots,2)\quad(p=3).
\]
Write \(q-7=\sum_j2r_jp^j\).  There is no carry on division by two, so
\[
 \sum_jr_jp^j=\frac{q-7}{2}=d-2.                       \tag{A.3}
\]
For \(2r_j\le p-1\), the submodule generated by
\((X^2)^{r_j}\) is the simple module
\(L(2r_j)=\Delta(2r_j)\), giving a Cartan embedding
\[
 L(2r_j)\lhook\joinrel\longrightarrow
 \operatorname{Sym}^{r_j}L(2).
\]
Tensoring its Frobenius twists and multiplying the polynomial factors
gives
\[
 L(q-7)\lhook\joinrel\longrightarrow
 \operatorname{Sym}^{d-2}L(2).                         \tag{A.4}
\]
This map is injective: in a monomial basis the three coordinate exponents
have base-\(p\) digits bounded by \(r_j<p\), so different tensor monomials
have different exponent triples.  Multiplication by the discriminant
\(Q\in\operatorname{Sym}^2L(2)\) is injective and completes
\[
 L(q-7)\lhook\joinrel\longrightarrow
 \operatorname{Sym}^{d-2}L(2)
 \mathrel{\mathop{\longrightarrow}^{\cdot Q}}F.        \tag{A.5}
\]

In the covariant \(\operatorname{GL}_2\)-convention, the source and target
degrees in (A.5) are \(q-7\) and \(q-3\), and \(Q\) transforms by
\(\det^2\).  Thus (A.5) is
\(\operatorname{GL}_2\)-equivariant from
\(L(q-7)\otimes\det^2\) to \(F\).  After the target is normalized by
\(\det^{-d}\), the source twist is
\[
 \det^{2-d}=\det^{-(q-7)/2}.                            \tag{A.6}
\]
Equation (A.1) shows that every finite-group map in this channel is
algebraic, so no occurrence can have the other outer normalization.
That other extension differs by \(\det^{(q-1)/2}\).

At \(q=9\), the same construction has a familiar model.  Under Hermite
reciprocity write \(A=X^3,B=X^2Y,C=XY^2,D=Y^3\).  The map is the
catalecticant-minor covariant
\[
 X^2\longmapsto B^2-AC,\qquad
 XY\longmapsto AD-BC,\qquad
 Y^2\longmapsto C^2-BD.
\]
It accounts for the \(L(2,0)=L(q-7)\) channel; no conclusion about the
remaining socle is drawn from this example.

For the prime-field statement, let \(q=p\ge5\).  Since \(d<p\), the
ordinary symmetrizer realizes \(F\) as a direct summand of
\(L(2)^{\otimes d}\).  Thus \(F\) is tilting.  The weight generating
function gives the Fischer character identity
\[
 \operatorname{ch}\operatorname{Sym}^dL(2)
 =\sum_{0\le j\le\lfloor d/2\rfloor}\chi_{2d-4j};
 \qquad
 \sum_{d\ge0}\operatorname{ch}(\operatorname{Sym}^dL(2))t^d
 =\frac1{(1-tz^2)(1-t)(1-tz^{-2})}.                    \tag{A.7}
\]
Every highest weight in (A.7) lies between \(0\) and \(p-3\), where the
indecomposable tilting module is simple.  Uniqueness of tilting
decomposition therefore yields
\[
 F=\bigoplus_{0\le j\le\lfloor(p-3)/4\rfloor}
 L(p-3-4j),                                             \tag{A.8}
\]
with multiplicity one; every summand has a projection.  This proves the
targeted linear-detector lemma.

\subsection{The opposite quadratic parity}
\label{app:outer-parity-proof}

We now treat maps into \(M=\operatorname{Sym}^2F\); the automatic
algebraicity argument above is not available.  Let \(S=L(s)\) be one of
the detector sources and put \(X=\operatorname{Hom}_k(S,M)\).
The polynomial degree of \(M\) is \(N=2q-6\).  With the covariant
normalization, a rational map from \(S\) uses determinant exponent
\((N-s)/2\); after tensoring source and target by their normalizing
determinant powers, the two projective extensions differ by
\(\det^{(q-1)/2}\).

The algebraic-torus weights in a \(\Gamma\)-fixed element of \(X\) are
multiples \(j(q-1)\), and
\[
 |\operatorname{wt}X|\le N+s\le3q-7<3(q-1).
\]
Thus a map \(\phi\) of the outer parity opposite to the normalized one has
a unique decomposition
\[
 \phi=\phi_++\phi_-,\qquad
 \operatorname{wt}(\phi_+)=q-1,\quad
 \operatorname{wt}(\phi_-)=1-q.                        \tag{A.9}
\]
For the positive root group define
\[
 D(t)=u_M(t)\phi-\phi u_S(t).
\]
Its entries have degree below \(2q\) and vanish on \(\F_q\), so
\[
 D(t)=(t^q-t)B(t),\qquad \deg B<q.                      \tag{A.10}
\]
The cocycle identity gives \(B(t+a)=B(t)u_S(a)\), first for
\(a\in\F_q\) and then for every \(a\in k\), since both sides have degree
at most \(q-1\) in \(a\).  Hence, for \(\beta_+=B(0)\),
\[
 u_M(t)\phi u_S(-t)-\phi=(t^q-t)\beta_+.               \tag{A.11}
\]
Comparing divided powers gives
\[
 E^{(1)}\phi_+=-\beta_+,\qquad E^{(q)}\phi_-=\beta_+,
 \qquad E^{(1)}\phi_-=E^{(q)}\phi_+=0.                  \tag{A.12}
\]

On the weight basis
\[
 A=x\otimes x^{(e)},\quad C=y\otimes x^{(e)},\quad
 B=x\otimes y^{(e)},\quad D=y\otimes y^{(e)}
\]
of \(Y_1\), choose the Weyl element with
\(wA=D,wC=-B,wB=-C,wD=A\), and define
\[
 \Theta(A)=-\beta_+,\quad \Theta(C)=\phi_+,\quad
 \Theta(B)=-\phi_-,\quad \Theta(D)=w\Theta(A).          \tag{A.13}
\]
The full divided-power table is
\[
\begin{array}{c|cccc}
 &\Theta(A)&\Theta(C)&\Theta(B)&\Theta(D)\\ \hline
 E^{(1)}&0&\Theta(A)&0&\Theta(B)\\
 E^{(q)}&0&0&\Theta(A)&\Theta(C)\\
 E^{(q+1)}&0&0&0&\Theta(A)\\ \hline
 F^{(1)}&\Theta(C)&0&\Theta(D)&0\\
 F^{(q)}&\Theta(B)&\Theta(D)&0&0\\
 F^{(q+1)}&\Theta(D)&0&0&0 .
\end{array}                                             \tag{A.14}
\]
Indeed, (A.12) gives the two middle columns and Weyl conjugation gives
their negative-root counterparts.  The ordinary divided-power commutator
gives \(F^{(1)}\Theta(A)=\Theta(C)\), since \(q-1=-1\) in \(k\).
Applying
\[
 E^{(q)}F^{(q)}
 =\sum_iF^{(q-i)}
   \binom{H-2q+2i}{i}E^{(q-i)}
\]
to \(\Theta(C)\) leaves \(i=q,q-1\), with coefficients \(0,1\), and
gives \(F^{(q)}\Theta(A)=\Theta(B)\).  Multiplication in the distribution
algebra gives \(F^{(q+1)}\Theta(A)=\Theta(D)\).

It remains to exclude unlisted divided powers.  Use the Chevalley
factorization
\[
 u(s)\ell(t)=
 \ell\!\left(\frac{t}{1+st}\right)h(1+st)
 u\!\left(\frac{s}{1+st}\right)                        \tag{A.15}
\]
when \(1+st\ne0\), then clear denominators.  Substitution of (A.12) and
the three corner coefficients above, followed by comparison of
coefficients of \(s^at^b\), gives
\[
 \ell(t)\Theta(A)=\Theta(A)+t\Theta(C)+t^q\Theta(B)
                    +t^{q+1}\Theta(D).
\]
Thus \(F^{(i)}\Theta(A)=0\) for
\(i\notin\{0,1,q,q+1\}\); Weyl conjugation and the middle-column
relations give every remaining zero in (A.14).  The weights, Weyl action,
and (A.14) prove that \(\Theta:Y_1\to X\) is a \(\mathbf G\)-map.
Moreover \(\Theta(C-B)=\phi\).  If \(\Theta=0\), then the root defect
vanishes and \(\phi\) is \(\mathbf G\)-fixed, impossible for the two
nonzero weights in (A.9).  We have therefore constructed the injection
\[
 \operatorname{Hom}_{\operatorname{PGL}_2(q)}(S^\square,M)
 \lhook\joinrel\longrightarrow
 \operatorname{Hom}_{\mathbf G}(S\otimes Y_1,M).        \tag{A.16}
\]

We verify the right side is zero source by source.  For
\(S=\operatorname{St}\), Doty--Henke's bottom-alcove tensor rule and
canonical tilting factorization give
\[
 \operatorname{St}\otimes L(1)\simeq T(q),\qquad
 \operatorname{St}\otimes Y_1\simeq T(2q).
\]
These are the bottom-alcove tensor identities of
\cite[Lemmas~1.3--1.4]{DotyHenke2002}.
The canonical tilting digits of \(2q\) are
\((p,p-1,\ldots,p-1,1)\).  The module \(T(p)\) has the two Weyl sections
\(\Delta(p)\) and \(\Delta(p-2)\)
\cite[Lemma~1.1]{DotyHenke2002}; hence \(T(2q)\) has exactly
\[
 \Delta(2q),\qquad \Delta(2q-2),                        \tag{A.17}
\]
once each.  Both highest weights exceed the top weight \(2q-6\) of \(M\).
Applying \(\operatorname{Hom}_{\mathbf G}(-,M)\) successively to this
two-section Weyl filtration gives
\[
 \operatorname{Hom}_{\mathbf G}(\operatorname{St}\otimes Y_1,M)=0.
\]

For \(q=p\) and \(S=L(s)\), \(s\in\{2,4,6,8,12\}\) in the ranges of
(3.2f), the only simple sources in \(S\otimes Y_1\) are
\(L(s+1,1)\) and \(L(s-1,1)\).  The module \(M\) is tilting:
\(F\) is tilting by (A.8), and the symmetric square is a direct summand
of the tensor square because \(p\) is odd.  Every indecomposable summand
\(T(n)\) of \(M\) has \(n\le2p-6\), and its restricted simple socle is
\(L(n)\) for \(n<p\), or \(L(2p-2-n)\) for \(n\ge p\).  Neither
two-digit candidate embeds, so the Hom space in (A.16) is zero.

Finally take \(S=L(q-7)\).  If \(p\ge7\), then
\[
 L(q-7)\otimes L(1)=L(q-6)\oplus L(q-8),
\]
with the second term omitted for \(p=7\); tensoring by \(L(1)^{(e)}\)
gives the sources \(L(2q-6)\) and \(L(2q-8)\).

We make the plethystic normalization explicit.  Write
\[
 e_i=X^{d-i}Y^i,\qquad
 a_{ij}=e_i\odot e_j\in\operatorname{Sym}^2D,\qquad
 [i,j\mid k,l]=a_{ij}\odot a_{kl}\in M,
\]
with the indices within each pair, and then the two pairs, ordered.
These are nested symmetric products with no hidden factor of \(2\);
in particular \([0,0\mid3,3]\) and \([0,3\mid0,3]\) are distinct.  Put
\[
 X_0=[0,0\mid0,0],\qquad Y_0=[0,0\mid0,1].
\]
For the negative and positive root groups,
\[
\begin{aligned}
 \ell(t)e_i&=\sum_{r=0}^{d-i}\binom{d-i}{r}t^re_{i+r},\\
 u(t)e_i&=\sum_{r=0}^{i}\binom{i}{r}t^re_{i-r};
\end{aligned}                                           \tag{A.17a}
\]
hence \(E^{(1)}Y_0=X_0\).

For \(p\ge7\), set
\[
 a=\frac{p-3}{2},\qquad b=\frac{p-7}{2}.
\]
The highest-weight space of \(M\) is \(kX_0\).  A highest vector \(g\) of
\(L(2q-6)\) satisfies \(F^{(p-5)}g=0\), because the least base-\(p\)
digit \(p-5\) of the divided-power index exceeds the least digit \(p-6\)
of its highest weight.  On the target, writing
\(A(t)=\ell(t)(e_0\odot e_0)\), the coefficient of
\([0,0\mid b,a]\) in \(F^{(p-5)}X_0\) is
\[
 4\binom db\binom da.                                  \tag{A.17b}
\]
Indeed, \(a+b=p-5\); the two factors of \(2\) are respectively the
inner and outer symmetric-square cross terms.  Since \(d\equiv a\pmod p\),
Lucas' theorem gives
\[
 \binom da=1,\qquad \binom db=\binom ab=\binom a2,
\]
so (A.17b) is nonzero, including when \(p=7\) and \(b=0\).
Thus the image of \(g\) must vanish.  The weight-\((2q-8)\) space is
\(kY_0\), but \(E^{(1)}Y_0=X_0\ne0\); hence no highest vector of
\(L(2q-8)\) can map nontrivially either.

If \(p=3\), write
\[
 L(q-7)\otimes Y_1=T(3)\otimes R^{(2)},
\]
where the previously suppressed higher-digit factor is
\[
 R^{(2)}
 =\left(\bigotimes_{j=2}^{e-1}L(2)^{(j)}\right)
   \otimes L(1)^{(e)};
\]
the digit \(j=1\) is trivial.  The module \(T(3)\) has Weyl sequence
\(0\to\Delta(3)\to T(3)\to\Delta(1)\to0\).  Let \(g\) be the highest generator
of the \(\Delta(3)\)-section, tensored with the highest vector of
\(R^{(2)}\), and let \(v\) lift the highest generator of the
\(\Delta(1)\)-quotient, normalized by \(E^{(1)}v=g\).  They have weights
\(2q-6,2q-8\) and generate the source.  The matching target spaces are
spanned by \(X_0,Y_0\), so a map has
\(g\mapsto\alpha X_0,\ v\mapsto\beta Y_0\).  But \(F^{(6)}g=0\), while
\[
 \operatorname{coeff}_{[0,0\mid3,3]}(F^{(6)}X_0)
 =2\binom d3^2=2.                                      \tag{A.17c}
\]
Here the factor \(2\) is the outer cross term; Lucas' theorem gives
\(\binom d3=1\) from the base-\(3\) digits of
\(d=(3^e-3)/2\).  Hence \(\alpha=0\), and
\(E^{(1)}v=g,\ E^{(1)}Y_0=X_0\) give \(\beta=0\).
This proves Lemma~\ref{lem:opposite-parity-quadratic}.

\subsection{Contraction and the prime-field affine class}
\label{app:contraction-proof}

Let \(i:S\hookrightarrow F\) have an \(H\)-retraction
\(\rho:F\to S\).  Define
\[
 D_\rho:\operatorname{Sym}^2F\longrightarrow S\otimes F,\qquad
 D_\rho(xy)=\rho(x)\otimes y+\rho(y)\otimes x.
\]
Compose \(D_\rho\) with a cochain
\(S\to\operatorname{Sym}^2F\), then take the categorical trace over
\(S\).  This gives an \(H\)-map
\[
 C_{\rho,i}:\operatorname{Hom}(S,\operatorname{Sym}^2F)\longrightarrow F.
\]
For the cocycle \(z_i(g)(t)=i(t)c(g)\) from (3.2b), the first summand
closes the \(S\)-loop and the second projects \(c(g)\).  Therefore
\[
 C_{\rho,i}(z_i)(g)=(\dim S)c(g)+i\rho(c(g)),           \tag{A.18}
\]
and the same identity holds in cohomology.

It remains to identify the prime-field class and its support.  Put
\[
 r=\frac{p-1}{2},\qquad m=r+1,\qquad
 P_0=Y(X^r-Z^r)\in R_m,\qquad R_m=k[X,Y,Z]_m.
\]
Let \(D=s^pt-st^p\) and
\[
 E_{\rm poly}=\{P\in R_m:P\circ\nu\in kD\}.
\]
With the determinant normalization of (2.2), restriction to the conic is
the exact \(G\)-sequence
\[
 0\longrightarrow QR_{m-2}\longrightarrow E_{\rm poly}
 \mathrel{\mathop{\longrightarrow}^{\epsilon_{\rm poly}}}k
 \longrightarrow0,                                    \tag{A.19}
\]
where \(\epsilon_{\rm poly}(P)=\lambda\) when
\(P\circ\nu=\lambda D\).  Since the conic ideal is \((Q)\),
multiplication by \(Q\) identifies the kernel with \(F=R_{m-2}\).
The affine chart \(\epsilon_{\rm poly}=1\) is the ambient point-vector
extension in (3.2d).

The base form \(P_0\) restricts to \(D\), and its connecting cocycle is
\[
 c(u(z))=\frac{u(z)P_0-P_0}{Q},\qquad
 u(z):(X,Y,Z)\mapsto(X,Y+zX,Z+2zY+z^2X).
\]
The coefficient of \(z^{p-2}\) is
\[
 [z^{p-2}]c(u(z))=-rX^{r-1}\ne0.                       \tag{A.20}
\]
Indeed the numerator coefficient is \(-rX^{r-1}Q\): the terms containing
one or two copies of \(2zY\) sum to
\(-2r^2X^{r-1}Y^2\), the term containing \(Z\) is \(-rX^rZ\), and
\(2r=-1\) in \(k\).  A coboundary \(u(z)A-A\), \(A\in R_{m-2}\), has
degree at most \(p-3\) in \(z\), so it cannot cancel (A.20).
Thus \([c]\ne0\).

Finally restrict cohomology to the prime-field Borel.  On \(L(n)\), the
root group \(U=C_p\) is one Jordan block, so \(H^1(U,L(n))\) is its
coinvariant line; the torus acts on that line by \(a^{-(n+2)}\).
For even \(0\le n\le p-3\), this line is fixed exactly when \(n=p-3\).
Restriction from \(H\) is injective because the Borel index is \(p+1\),
which is nonzero in characteristic \(p\).  Hence \([c]\) is supported
only on the top Fischer summand in (A.8).  Formula (A.18) is therefore
\((s+1)[c]\) on a lower detector and \((s+2)[c]\) on the top detector.
The ranges in (3.2f) make the first scalar nonzero; in the top case
\(s=p-3\), so the second is \(p-1\ne0\).  This proves
Lemma~\ref{lem:affine-class-contraction}.

\section{Six profiles and matching-row reconstruction}
\label{sec:six-profiles}

Specialize to \(H_3\).  Put
\[
 G=\operatorname{PGL}_2(11),\qquad
 G^+=\operatorname{PSL}_2(11),\qquad
 X=\Omega_{H_3}\simeq G/H,
\]
where \(H\simeq A_5\) is the stabilizer of the base matching in
(3.1).  Thus \(X=X_+\sqcup X_-\), with \(|X_+|=|X_-|=11\).
The construction has four stages:
\[
 (A_+,A_-)\longmapsto K
 \longmapsto\text{four oriented cell pairs}
 \longmapsto D(M)
 \longmapsto K\backslash G/H.
\]
The ordered parent pair supplies the \(A_4\)-subgroup \(K\); its cell
partition supplies four signed incidence counts; those counts recover
the double-coset label.  Only the first two stages require coordinates.

\subsection{The decorated parent correspondence}

On
\(V=\F_{11}^3\), set
\[
 q_0(x,y,z)=x^2+y^2+z^2,\qquad
 \cC_0=\mathbb P\{q_0=0\}.
\]
The projectivity
\[
 \psi=
 \begin{pmatrix}10&2&1\\1&2&10\\3&0&3\end{pmatrix}
 \quad\text{satisfies}\quad
 q_0(\psi(X,Y,Z)^t)=3(XZ-Y^2),                         \tag{B.1}
\]
so \(\psi\circ\nu\) identifies the endpoint conic of
Section~\ref{sec:matching-products} with
\(\cC_0\).  In these coordinates take the two six-point sets
\[
\begin{aligned}
A_+={}&\{(0:1:4),(0:1:7),(1:0:3),(1:0:8),
          (1:4:0),(1:7:0)\},\\
A_-={}&\{(0:1:3),(0:1:8),(1:0:4),(1:0:7),
          (1:3:0),(1:8:0)\}.
\end{aligned}                                           \tag{B.2}
\]
Let \(\mathcal P=G\cdot A_+\), where \(G\) is viewed as the full
projective stabilizer of \(\cC_0\).  For \(A\in\mathcal P\) and
\(a\in A\), let \(C_{A,a}\) be the conic through
\(A\setminus\{a\}\).  The six pairs
\[
 C_{A,a}\cap\cC_0,\qquad a\in A,                        \tag{B.3}
\]
form a perfect matching \(M_A\) of the twelve points of \(\cC_0\).
On this \(22\)-point orbit, \(A\mapsto M_A\) is a
\(G\)-equivariant bijection onto \(X\).  This is the
\emph{matching-decorated parent correspondence}; its bijectivity and
the assertion that (B.3) consists of six disjoint pairs are exact
finite inputs.  Concretely, the supplement enumerates the \(22\) translates
of \(A_+\), deletes each of their six points, solves the resulting
five-point conic, intersects it with \(\cC_0\), rejects unless the six
intersections are disjoint two-sets, and compares the resulting canonical
matching with all \(22\) rows of \(X\).  It is the secant-factorization counterpart of the
self-polar matching and Brianchon-support reconstruction obtained from
the Clebsch decoding geometry \cite{RuddRigidity2026}.

Put \(H_\pm=\operatorname{Stab}_G(A_\pm)\) and
\[
 K=H_+\cap H_-,\qquad
 J=\begin{pmatrix}1&0&0\\0&0&-1\\0&-1&0\end{pmatrix}.
                                                               \tag{B.4}
\]
Then \(K\simeq A_4\), \(J(A_+)=A_-\), and \(J\) normalizes \(K\) and
exchanges \(X_+\) with \(X_-\).  Thus the \emph{ordered golden pair}
means precisely the ordered pair \((A_+,A_-)\).

Choose linear representatives of \(K\) preserving \(q_0\), and let
\(\widetilde K=\F_{11}^{\times}K\).  For the following eight vectors,
define \(R_i=\mathbb P(\widetilde K r_i)\):
\[
\begin{array}{c@{\quad}c@{\qquad}c@{\quad}c}
i&r_i&i&r_i\\ \hline
1&(0,1,4)&10&(0,1,3)\\
3&(0,1,2)&13&(0,1,5)\\
6&(1,3,2)&14&(1,1,5)\\
9&(1,4,2)&11&(1,2,4).
\end{array}                                               \tag{B.5}
\]
These are cells in the scalar-\(K\) orbit partition of
\(\mathbb P(V)\).  The involution \(J\) exchanges the four pairs
\[
 (R_1,R_{10}),\quad (R_3,R_{13}),\quad
 (R_6,R_{14}),\quad (R_9,R_{11}).                         \tag{B.6}
\]
Only these four oriented pairs enter the argument.  Reversing the
golden pair reverses all four orientations.

This profile layer uses more decoration than the balanced-sheet theorem.
The quotient configuration intrinsically recovers the unordered sheets,
but it does not select an ordered golden pair or the resulting
scalar-\(K\) refinement.  The reconstruction below is relative to that
selected refinement.

\subsection{Incidence profiles}

For \(M\in X\), let \(S_M\) be the union of its six secants in
\(\mathbb P^2(\F_{11})\), and put
\[
 Z_M(R)=|S_M\cap R|.
\]
Equivalently, \(Z_M(R)\) counts the points of \(R\) at which \(P_M\)
vanishes.  Define the \emph{depth profile}
\[
 \begin{aligned}
 D(M)=\bigl(&Z_M(R_1)-Z_M(R_{10}),\
              Z_M(R_3)-Z_M(R_{13}),\\
            &Z_M(R_6)-Z_M(R_{14}),\
              Z_M(R_9)-Z_M(R_{11})\bigr)\in\F_{11}^4.
 \end{aligned}                                           \tag{B.7}
\]
The entries here and in the profile table below are first formed as
integer differences of cardinalities, using the canonical cell
orientation in (B.6), and only then reduced modulo \(11\).
This definition uses only zero sets of secant products, so it is
unchanged if their equations are rescaled.

\begin{theorem}[Six-profile reconstruction]
\label{thm:six-profile-reconstruction}
The \(K\)-orbits on each sheet \(X_\pm\) have sizes \(1,4,6\).
With the orientation in (B.6), the corresponding profiles and fibre
sizes are
\[
\begin{array}{c@{\quad}c@{\quad}c}
\text{sheet}& |D^{-1}(v)|&v\\ \hline
+&1&v_1=(-6,0,12,-12),\\
+&4&v_2=(-3,3,0,3),\\
+&6&v_3=(3,-2,-2,0),\\
-&1&-v_1,\\
-&4&-v_2,\\
-&6&-v_3.
\end{array}                                                \tag{B.8}
\]
The six vectors are pairwise distinct.  Hence \(D\) recovers the
double-coset row in
\[
 K\backslash G/H
\]
of every matching.  Its linear span nevertheless has dimension only
two: over \(\F_{11}\) it is the plane
\[
 2a+2b+c=0,\qquad 9a+8b+d=0.                              \tag{B.9}
\]
The two singleton profiles recover their individual matchings.  Under
the matching-decorated parent correspondence, they therefore recover
the corresponding individual \(H_3\) parent.  A profile in a fibre of
size four or six recovers only that \(K\)-orbit, not an individual
matching.
\end{theorem}

\noindent\emph{Computer-assisted proof mode.}
The subgroup-mark reduction and the implications from equivariance are
conceptual.  The six representative incidence rows in (B.10) below are
certificate-checked exact counts and are reconstructed by an independent
implementation.  A companion certificate verifies that the matching
decoration uniquely recovers its parent on the \(22\)-point \(H_3\)
locus.

\begin{proof}
The permutation character of \(K\) on either eleven-point sheet has
values
\[
 (11,3,2)
\]
on elements of orders \(1,2,3\).  The marks of the transitive
\(A_4\)-sets are
\[
\begin{array}{c@{\quad}c@{\quad}c}
\text{\(K\)-set}&\text{size}&\text{marks on orders \(1,2,3\)}\\ \hline
K/K&1&(1,1,1)\\
K/V_4&3&(3,3,0)\\
K/C_3&4&(4,0,1)\\
K/C_2&6&(6,2,0)\\
K/1&12&(12,0,0).
\end{array}
\]
Comparing marks gives the unique nonnegative decomposition
\[
 X_+|_K\simeq K/K\ \sqcup\ K/C_3\ \sqcup\ K/C_2.
\]
The orbit sizes are therefore \(1,4,6\); conjugation by \(J\) gives the
same decomposition on \(X_-\).

For \(k\in K\), the transformation \(k\) preserves every cell in
(B.6) and carries \(S_M\) to \(S_{kM}\).  Thus \(D(kM)=D(M)\).
The involution \(J\) exchanges the two cells in every oriented pair and
carries \(S_M\) to \(S_{JM}\), so
\[
 D(JM)=-D(M).
\]
It remains to count one representative on each positive-sheet orbit.
The exact incidence rows are
\[
\begin{array}{c@{\quad}c@{\quad}c@{\quad}c}
\text{orbit size}&\operatorname{Stab}_K(M)&
 \text{zero counts in the four pairs of (B.6)}&D(M)\\ \hline
1&A_4&(0,6;\ 0,0;\ 12,0;\ 0,12)&v_1,\\
4&C_3&(3,6;\ 3,0;\ 3,3;\ 6,3)&v_2,\\
6&C_2&(3,0;\ 1,3;\ 2,4;\ 6,6)&v_3.
\end{array}                                               \tag{B.10}
\]
The \(J\)-images supply the three negative rows.  These six vectors are
pairwise distinct, so their fibres are precisely the six \(K\)-orbits.

Row reduction over \(\F_{11}\) gives rank two and the equations (B.9).
This rank drop does not merge orbit labels: linear rank and
set-theoretic separation are different assertions.  Finally, a
singleton fibre contains its matching without residual choice.  The
stabilizer of the matching \(M_A\) from (B.3) is the same \(A_5\) as
the stabilizer of its parent \(A\).  The equivariant decorated transform between
the two \(G/H\)-orbits is therefore a bijection, so the matching also
determines its parent.  The stated limitation for the other four rows
follows from their fibre sizes.
\end{proof}

\begin{corollary}[Decorated sheet classifier]
\label{cor:decorated-sheet-classifier}
Relative to the orientation in (B.6), the profile \(D(M)\) determines
the sheet containing \(M\).  The antipodal singleton pair
\(\{v_1,-v_1\}\) recovers the unordered golden matching pair and hence
the unordered golden parent pair; choosing the orientation selects one
member.
\end{corollary}

\begin{proof}
The positive and negative profile sets in (B.8) are disjoint, and the
only singleton fibres are \(v_1\) and \(-v_1\).  Apply the
matching--parent bijection from
Theorem~\ref{thm:six-profile-reconstruction}.
\end{proof}

Although the ambient radial coordinate separates the sheets, \(D\) uses
additional scalar-\(K\) incidence data and is not an affine-linear functional
on the abstract configuration.  The quadratic and cubic moments remain the
intrinsic recovery mechanism of Section~\ref{sec:balanced}.

The profile compression also retains the cubic-first orientation from
Section~\ref{sec:balanced}.  The three positive vectors satisfy the
integral weighted relation
\[
 v_1+4v_2+6v_3=0.                                       \tag{B.11}
\]
\begin{corollary}[Orbit weights from normalized profiles]
\label{cor:profile-ray-weights}
The three incidence-normalized vectors \(v_1,v_2,v_3\) recover the
orbit-size multiset \(\{1,4,6\}\).  They also recover the stabilizer
orders \(\{12,3,2\}\) inside \(K\simeq A_4\).  The underlying rational
rays alone do not recover these weights.
\end{corollary}

\begin{proof}
The coordinates of each \(v_i\) are fixed by the cardinality
differences in (B.7), so independent rescaling is not allowed.  These
three integral vectors span a plane over \(\mathbb Q\), and (B.11) is
the primitive positive generator of their relation lattice.  Its
coefficients recover the three orbit sizes without labels.
Orbit--stabilizer then gives \(12/1,12/4,12/6\).  If the vectors are
replaced by primitive generators of their rays, the relation becomes
\(1:2:1\), which proves the final warning.
\end{proof}

\section{Modular depth and the projective-cover bridge}
\label{sec:modular-depth}

Because an \(H_3\) sheet has eleven points in characteristic eleven,
the sum of its basis vectors is simultaneously constant and
augmentation-zero.  Linearizing the three incidence rows therefore
kills exactly this constant direction and leaves the two-dimensional
profile plane.  The following module statement identifies that loss.

Retain the notation of Section~\ref{sec:six-profiles}, put
\[
 k=\F_{11},\qquad L=G^+=\operatorname{PSL}_2(11),
 \qquad V=k[X_+]\simeq\operatorname{Ind}_H^L\mathbf 1,
\]
and regard $V$ as the permutation module on one eleven-point sheet.
Let $e_1,e_4,e_6$ be the indicator functions of the three $K$-orbits,
in orbit-size order.  Thus
\[
 V^K=ke_1\oplus ke_4\oplus ke_6,\qquad
 \mathbf 1=e_1+e_4+e_6.                                \tag{C.1}
\]
Linearizing the incidence construction in (B.7) gives
\[
 \mathcal D:V^K\longrightarrow k^4,\qquad
 e_i\longmapsto i\,u_i,                                \tag{C.2}
\]
where $u_1=v_1$, $u_4=v_2$, and $u_6=v_3$ are the profiles on the
orbits of the indicated sizes.
Equation (B.11) says exactly that $\mathcal D(\mathbf1)=0$.

\begin{proposition}[Modular depth quotient]
\label{prop:modular-depth-quotient}
The $L$-module $V$ is the projective cover $P(\mathbf1)$ and has
Loewy dimensions
\[
 1\mid9\mid1.
\]
The map $\mathcal D$ has kernel $k\mathbf1$ and induces an isomorphism
\[
 P(\mathbf1)^K/\operatorname{soc}P(\mathbf1)
 \ \xrightarrow{\ \sim\ }\ \langle u_1,u_4,u_6\rangle_k. \tag{C.3}
\]
The sheet exchanged by $J$ carries the negative copy of this quotient.
\end{proposition}

\noindent\emph{Proof mode.}
Projectivity, localness, and fixed-point exactness are conceptual.  The
middle layer follows from the earlier \(M_8\)-identification.  The rank and
kernel of \(\mathcal D\) use the certified profile table (B.8) and its
independent replay.

\begin{proof}
The order of $H\simeq A_5$ is prime to $11$, so its trivial module is
projective and induction makes $V$ projective.  The action of $L$ on
$X_+$ is doubly transitive.  Hence $\operatorname{End}_L(V)$ is
generated by the identity and the all-ones operator $E$.  In
characteristic eleven,
\[
 E^2=11E=0,
\]
so this endomorphism algebra is local.  Thus $V$ is indecomposable.  Its
trivial head occurs once, and consequently $V=P(\mathbf1)$.  The
constant vector is its one-dimensional socle.  The unique head map is
augmentation, so
\[
 \operatorname{rad}V=\ker\!\left(V\xrightarrow{\sum}k\right).
\]
On this augmentation module define
\[
 \Phi\!\left(\sum_{M\in X_+}a_M[M]\right)
   =\sum_{M\in X_+}a_M\,\pi_8(x_M).
\]
Changing the reference matching translates every \(x_M\) by the same
vector, and the affine cocycle term cancels for the same reason; thus
\(\Phi\) is a well-defined \(G^+\)-map.  The same-sheet first-moment
identity of Proposition~\ref{prop:radical-hadamard} kills the constant
vector and shows that the induced map
\[
 \overline\Phi:\operatorname{rad}V/k\mathbf1\longrightarrow M_8
\]
is nonzero.  The target is the absolutely irreducible \(G^+\)-module
identified in Section~\ref{sec:rank-three}; both sides have dimension
nine.  Hence \(\overline\Phi\) is an isomorphism.  This proves the
displayed Loewy dimensions without a generator-matrix calculation.

Taking $K$-fixed points is exact because $11\nmid |K|=12$.  The orbit
sizes $1,4,6$ give (C.1).  The three vectors in (C.2) span the plane
 (B.9), and their only relation is (B.11), so
$\ker\mathcal D=k\mathbf1$.  This proves (C.3).  Finally
$D(JM)=-D(M)$, which attaches the outer character on the opposite
sheet.
\end{proof}

\begin{corollary}[The canonical nine-space bridge]
\label{cor:h3-nine-space-bridge}
For the fixed marked conic and sheet \(X_+\), there is a canonical
\(G^+\)-module isomorphism
\[
 \operatorname{rad}P(\mathbf1)/\operatorname{soc}P(\mathbf1)
 \ \xrightarrow{\ \sim\ }\ M_8\!\downarrow_{G^+}.             \tag{C.4}
\]
It sends the augmentation class of
\(\sum_{M\in X_+}a_M[M]\) to
\[
 \sum_{M\in X_+}a_M\,\pi_8(x_M),
 \qquad \sum_Ma_M=0,                                         \tag{C.5}
\]
where \(\pi_8\) is the top-harmonic projection.  Formula (C.5) is
independent of the reference matching.
\end{corollary}

\begin{proof}
The construction and all asserted properties of (C.5) are precisely the
map \(\overline\Phi\) established in the proof of
Proposition~\ref{prop:modular-depth-quotient}.
\end{proof}

\begin{corollary}[The homogeneous projective-cover bridge]
\label{cor:h3-homogeneous-projective-cover}
For the cocycle representative \(c_8\) determined by the chosen
reference matching, let \(E_{c_8}\) be the \(G\)-module defined in
(3.11b).
There is a \(G^+\)-module isomorphism
\[
 P(\mathbf1)/\operatorname{soc}P(\mathbf1)
 \ \xrightarrow{\ \sim\ }\ E_{c_8}\!\downarrow_{G^+},          \tag{C.6}
\]
given by
\[
 \Theta\!\left(\sum_{M\in X_+}a_M[M]\right)
   =\left(\sum_{M\in X_+}a_M\pi_8(x_M),\ \sum_{M\in X_+}a_M\right).
                                                               \tag{C.7}
\]
Thus the extension class, unlike a cocycle representative or an affine
origin, is canonical: the top-harmonic homogeneous extension restricts
to the nonsplit exact sequence
\[
 0\longrightarrow M_8\!\downarrow_{G^+}
 \longrightarrow P(\mathbf1)/\operatorname{soc}P(\mathbf1)
 \longrightarrow \mathbf1\longrightarrow0 .
\]
If \(\widetilde E_c\) denotes the homogenization of the full affine
quotient with translation space \(M_8\oplus M_0\), then
\[
 \widetilde E_c\!\downarrow_{G^+}\simeq
 P(\mathbf1)/\operatorname{soc}P(\mathbf1)\oplus\mathbf1,      \tag{C.8}
\]
where the last summand is the radial module \(M_0\!\downarrow_{G^+}\).
\end{corollary}

\begin{proof}
The affine transformation law
\(\pi_8(x_{gM})=g\pi_8(x_M)+c_8(g)\) makes \(\Theta\)
\(G^+\)-equivariant.  The constant vector maps to zero: its second
coordinate is \(|X_+|=11=0\) in \(k\), and its first coordinate
vanishes by the same-sheet first-moment identity.  Hence \(\Theta\)
descends through the socle.

On the kernel of augmentation, the descended map is the isomorphism
of Corollary~\ref{cor:h3-nine-space-bridge}; on the one-dimensional
head it induces the identity through the second coordinate.  It is
therefore an isomorphism.  The exact sequence cannot split:
Proposition~\ref{prop:modular-depth-quotient} gives
\(P(\mathbf1)/\operatorname{soc}P(\mathbf1)\) the simple socle \(M_8\),
whereas a splitting would put a trivial summand in its socle.  Finally,
\(M_0=\chi\) is trivial on \(G^+\), and its cocycle vanishes there,
proving (C.8).
\end{proof}

\begin{remark}[Ambient and intrinsic sheet recovery]
With the harmonic--radial decomposition retained, the \(M_0\)-projection
already separates the two \(H_3\) sheets: it is zero on \(X_+\), constant
on \(X_-\), and nonzero by Lemma~\ref{lem:shared-radial-cycle}.  The quadratic theorem of
Section~\ref{sec:balanced} is stronger in a different sense: it recovers
the sheets intrinsically from the abstract affine point configuration,
without retaining the ambient Fischer operator.
\end{remark}

This quotient identifies the precise modular loss in the six-profile
construction.  Before linearization, the six profiles remain distinct
and recover all six double-coset labels.  After linearization, the
constant orbit trade dies and only the two-dimensional depth plane
remains.  Conversely, the three rational rays in that plane remain
distinct but do not retain their orbit weights.  The weights \(1,4,6\) are
recovered only when the quotient is supplied with the
incidence-normalized vectors of (B.8), as in
Corollary~\ref{cor:profile-ray-weights}.

\begin{remark}[Reduction of the profile table]
\label{rem:bad-profile-primes}
The four-coordinate profiles in (B.8) are integral before reduction.
The greatest common divisor of their nonzero $2$-by-$2$ minors is $3$.
Their span therefore drops below rank two only in characteristic $3$.
The six signed rows remain distinct away from characteristics $2$ and
$3$: characteristic $2$ identifies antipodes, while characteristic
$3$ sends the rows $v_1$ and $v_2$ to zero.  This is a statement about
the displayed integral table, not a construction of the marked conic
geometry in those characteristics.
\end{remark}

\section{Arithmetic splitting of the marker algebra}
\label{sec:arithmetic-splitting}

The arithmetic representatives of the three configurations come in
conjugate pairs.  Besides the base matchings in (3.1), we use
\[
\begin{aligned}
 M_{B,-}&=\{\{0,\infty\},\{1,3\},\{2,6\},\{4,5\}\},\\
 M_{B,+}&=\{\{0,\infty\},\{1,5\},\{2,3\},\{4,6\}\},\\
 M_{H,-}&=\{\{0,10\},\{1,\infty\},\{2,7\},\{3,5\},
                   \{4,8\},\{6,9\}\}.
\end{aligned}                                             \tag{D.1}
\]
The pair for $A_3$ consists of two copies of $M_{A_3}$, the pair for
$B_3$ is $(M_{B,-},M_{B,+})$, and the pair for $H_3$ is
$(M_{H_3},M_{H,-})$.  Put
\[
 f_{A_3}=f_{B_3}=t^2-2,\qquad f_{H_3}=t^2-t-1;
\]
the quadratic marker algebra of type \(T\) is \(\F_q[t]/(f_T)\).
To make the arithmetic bridge part of the definition, let \(\bar k\) be
an algebraic closure of the working field
and define the root-specialization maps
\[
\begin{array}{c@{\quad}c@{\quad}c}
T&\text{root}&\sigma_T(\text{root})\\ \hline
A_3&\sqrt2,\ -\sqrt2&M_{A_3},\ M_{A_3}\\
B_3&3,\ 4&M_{B,-},\ M_{B,+}\\
H_3&4,\ 8&M_{H_3},\ M_{H,-}.
\end{array}                                               \tag{D.2}
\]
Thus the \emph{marker datum} is the quadratic algebra together with
\(\sigma_T\), not the polynomial alone.  Exact projective transport
places the displayed $B_3$ and $H_3$ values in
\(\Omega_{B_3}\) and \(\Omega_{H_3}\).  The certificate also verifies
that \(\sigma_T\) is the reduction of the corresponding integral
Coxeter frames; the theorem below needs only the explicit map (D.2).
We call the datum \emph{fused} when its two geometric roots have one
rational matching value, and \emph{split} when they give two distinct
rational matching sheets.

\begin{lemma}[Split--inert frame calculation]
\label{lem:split-inert-frames}
The three marker polynomials factor as follows:
\[
\begin{array}{c@{\qquad}c@{\qquad}c}
A_3,\ q=5&t^2-2&\text{irreducible},\\
B_3,\ q=7&t^2-2=(t-3)(t-4)&\text{split},\\
H_3,\ q=11&t^2-t-1=(t-4)(t-8)&\text{split}.
\end{array}                                               \tag{D.3}
\]
In the first row the two geometric roots form one Frobenius orbit over
$\F_{25}$.  In each split row the nontrivial involution of the
quadratic marker algebra exchanges the two rational roots.
\end{lemma}

\begin{proof}
The nonzero squares modulo $5$ are $1$ and $4$, so $2$ is not a square.
If $u^2=2$ in $\F_{25}$, the nontrivial $\F_5$-automorphism sends $u$
to $-u$, and hence exchanges the two roots.  The remaining
factorizations follow by direct substitution.
\end{proof}

\begin{theorem}[Marker specialization and \(H_3\) arithmetic gluing]
\label{thm:rank-three-arithmetic-gluing}
For these arithmetic representatives of the rank-three matching
configurations, the following hold.
\begin{enumerate}[label=\textup{(\alph*)},leftmargin=*]
\item The $A_3$ marker datum is fused over $\F_5$.  The $B_3$ and
$H_3$ marker data split into two rational matching sheets over $\F_7$
and $\F_{11}$, respectively.
\item In the $H_3$ case, the stabilizers $H_+$ and $H_-$ of the
singleton matching in each rational sheet are isomorphic to $A_5$,
satisfy
\[
 H_+\cap H_-=K\simeq A_4,\qquad
 \langle H_+,H_-\rangle=\operatorname{PSL}_2(11),
                                                               \tag{D.4}
\]
and are exchanged by a projectivity $J$ of nonsquare determinant.  The
normalizer \(N_G(K)\) is a rational \(S_4\) whose determinant-square
kernel is \(K\).
\item The determinant character that distinguishes the two rational
sheets also negates the depth profiles.  Thus choosing one root in the
split marker algebra chooses one sheet, while its singleton profile
chooses the corresponding decorated matching row.  Forgetting that
choice leaves the unordered pair.
\end{enumerate}
\end{theorem}

\noindent\emph{Computer-assisted proof mode.}
Lemma~\ref{lem:split-inert-frames} and the determinant-character
deductions are conceptual.  The agreement of the specialization map
(D.2) with the integral Coxeter frames, the matching stabilizers and
intersections, and the generated subgroup in (D.4) are exhaustive
exact finite checks.
The classical subgroup names use the standard subgroup structure of
$\operatorname{PGL}_2(11)$ \cite{Atlas1985,Giudici2007}.

\begin{proof}
Part (a) records the specialization data in (D.2).  Its two
\(A_3\) roots have the same matching value.  The two values in the
\(B_3\) row are exchanged by
\(x\mapsto-x\), and the two values in the \(H_3\) row are exchanged by
\[
 x\longmapsto\frac{x-1}{x+1}.
\]
The exact transport check verifies these identities on the displayed
matchings.  Together with Lemma~\ref{lem:split-inert-frames}, they prove
(a).

For $H_3$, exact projective closure gives
\[
 |H_+|=|H_-|=60,\qquad |H_+\cap H_-|=12,\qquad
 |\langle H_+,H_-\rangle|=660.
\]
Both stabilizers lie in the determinant-square subgroup, while the
displayed transporter has nonsquare determinant.  The normalizer
\(N_G(K)\) has order $24$, and its determinant-square kernel has order
$12$.  The cited subgroup identifications give (D.4) and the
$S_4/A_4$ hinge, proving (b).

The determinant-square subgroup has the two matching sheets as its
orbits, and $J$ exchanges them.  Section~\ref{sec:six-profiles} proves
$D(JM)=-D(M)$ and identifies the unique singleton in each sheet.
Theorem~\ref{thm:six-profile-reconstruction} then returns the decorated
matching row from that singleton.  This proves (c).
\end{proof}

Arithmetic splitting and modular depth retain different amounts of
information.  Splitting produces a two-point choice of sheet.  The
depth quotient is defined only after the $A_4$-refinement has been
selected and remembers all three orbit rays modulo the constant socle.
The first choice orients the second, but does not canonically identify
it with any cubic quotient of the ten-dimensional factorization space.
Appendix~\ref{app:tate-plane} makes that boundary precise.

\section{Cubic structure of the profile plane}
\label{app:profile-cubic}

\begin{lemma}[Three-ray cubic forcing]
\label{lem:three-ray-cubic}
Let \(k\) have characteristic different from \(2\) and \(3\).  Suppose
\(u_1,u_2\) form a basis of a two-dimensional \(k\)-space and
\(n_1,n_2,n_3\in k^\times\) satisfy
\[
 n_1u_1+n_2u_2+n_3u_3=0.
\]
For the signed antipodal measure
\[
 \xi=\sum_{i=1}^3 n_i(\delta_{u_i}-\delta_{-u_i}),
\]
the moments of degrees one and two vanish, while the cubic moment is
nonzero.
\end{lemma}

\begin{proof}
The relation kills the first moment, and antipodality kills every even
moment.  Write
\[
 u_3=-\frac{n_1}{n_3}u_1-\frac{n_2}{n_3}u_2.
\]
In \(\sum_i n_i u_i^{\odot3}\), the coefficient of
\(u_1^{\odot2}\odot u_2\) is
\[
 -\frac{3n_1^2n_2}{n_3^2},
\]
which is nonzero.  The signed cubic is twice this tensor and is therefore
nonzero as well.
\end{proof}

\begin{corollary}[The mass-zero cubic flag]
\label{cor:mass-zero-cubic}
In Lemma~\ref{lem:three-ray-cubic}, suppose also that
\(n_1+n_2+n_3=0\).  Then its half-cubic factors as
\[
 \frac12M_3(\xi)=\sum_{i=1}^3n_i u_i^{\odot3}
 =\frac{n_1n_2}{(n_1+n_2)^2}
 (u_1-u_2)^{\odot2}\odot
 \bigl((2n_1+n_2)u_1+(n_1+2n_2)u_2\bigr).               \tag{E.1}
\]
The two displayed lines are distinct.  Hence the cubic intrinsically
defines a doubled line and a residual simple line; its Hessian recovers
the doubled line.
\end{corollary}

\begin{proof}
Since \(n_3=-(n_1+n_2)\), the weighted vector relation gives
\[
 u_3=\frac{n_1u_1+n_2u_2}{n_1+n_2}.
\]
Substitution and expansion give (E.1).  The residual factor would equal
the doubled factor only if \(3(n_1+n_2)=0\), which is impossible under
the hypotheses.  The Hessian of a binary cubic \(L^2R\), with \(L\) and
\(R\) distinct, is a nonzero scalar multiple of \(L^2\).
\end{proof}

For the incidence-normalized profiles of Section~\ref{sec:six-profiles},
put
\[
 \eta=\delta_{v_1}+4\delta_{v_2}+6\delta_{v_3}
      -\delta_{-v_1}-4\delta_{-v_2}-6\delta_{-v_3}.
\]
Then
\[
 M_j(\eta)=
 \bigl(1-(-1)^j\bigr)
 \sum_{i=1}^3 n_i v_i^{\odot j},
 \qquad (n_1,n_2,n_3)=(1,4,6).                           \tag{E.2}
\]
Equation (B.11) kills the first moment, and antipodality kills every even
moment.  The first two profiles are independent over \(\F_{11}\), since
the determinant of their first two coordinates is \(-18\ne0\).
Lemma~\ref{lem:three-ray-cubic} therefore forces the cubic to be nonzero.
In the frozen normalization, its first coordinate is
\[
2\bigl((-6)^3+4(-3)^3+6(3)^3\bigr)=6\ne0
\quad\text{in }\F_{11}.                                  \tag{E.3}
\]
Thus the six-row quotient preserves the cubic-first signed moment even
though its linear image is only a plane.

\begin{remark}[The \(H_3\) cubic flag]
\label{rem:cubic-target-flag}
Here \(1+4+6=0\) in \(\F_{11}\), so
Corollary~\ref{cor:mass-zero-cubic} forces the double-line factorization.
In the basis \(e_1=v_2,e_2=v_3\) of the profile plane, the projective
cubic is
\[
 f(x,y)=x^3+7x^2y+5xy^2+9y^3=(x-y)^2(x-2y),
\]
with Hessian \(7x^2+8xy+7y^2=7(x-y)^2\).  Hence the cubic
intrinsically distinguishes a doubled line and a residual simple line.
This flag is internal to the decorated profile plane: it supplies no
canonical map from the full relative-cubic space and no descended
\(G^+\)-action.
\end{remark}

\section{The relative-cubic Tate plane}
\label{app:tate-plane}

Let $W=W_{H_3}$ be the ten-dimensional quotient space from
Theorem~\ref{thm:rank-three-quotients}, let $\chi$ be the determinant
character of $G=\operatorname{PGL}_2(11)$, and put
\[
 M=\operatorname{Sym}^3(W)\otimes\chi,\qquad R=M^G.
\]
The signed cubic of Section~\ref{sec:balanced} lies in $R$, which has
dimension three.  There is a canonical two-dimensional quotient of
this source space, but it is not canonically the depth plane of
Proposition~\ref{prop:modular-depth-quotient}.

We fix the monomial basis of \(\operatorname{Sym}^3(W)\) induced by
the harmonic--radial coordinates of Section~\ref{sec:rank-three}.
The \emph{frozen relative basis} is the reduced-echelon basis of the
simultaneous equations \(gf=\chi(g)f\) for the standard generators
defining the marked \(H_3\) orbit, ordered by its pivots.  Thus the
coordinates below specify a reproducible basis, not an additional
intrinsic structure.  The \emph{three labelled cubic orbit classes}
are indexed, in order, by the three \(K\)-orbits of sizes \(1,4,6\)
on \(X_+\).  Finally, the \emph{invariant symmetric-form pencil} is
\(\mathbb P((\operatorname{Sym}^2W)^G)\).  Its unique rank-one point
has the form \([\ell^2]\); contraction by its covector \(\ell\) is
the induced map
\[
 \iota_\ell:R\longrightarrow(\operatorname{Sym}^2W)^G.
\]

\begin{proposition}[Canonical relative-cubic Tate plane]
\label{prop:relative-cubic-tate-plane}
The invariant space $R$ and the coinvariant space $M_G$ both have
dimension three.  The canonical projection and group norm fit into
\[
 R\xrightarrow{\ \pi\ }M_G\xrightarrow{\ N\ }R,
\qquad
\operatorname{rank}\pi=2,\quad\operatorname{rank}N=1,
                                                               \tag{F.1}
\]
with
\[
 \operatorname{im}\pi=\ker N,\qquad
 \operatorname{im}N=\ker\pi=\langle[1:3:9]\rangle.             \tag{F.2}
\]
Contraction by the covector whose square is the unique rank-one member
of the invariant symmetric-form pencil gives a second construction of
the same quotient.  In the frozen relative basis its matrix is
\[
 \begin{pmatrix}8&1&0\\0&8&1\end{pmatrix},
                                                               \tag{F.3}
\]
and its kernel is again $\langle[1:3:9]\rangle$.
\end{proposition}

\noindent\emph{Computer-assisted proof mode.}
The invariant/coinvariant construction and the implications of the
displayed ranks are conceptual.  The dimensions, norm and projection
ranks, kernel line, rank-one pencil member, contraction matrix, and
flattening census are exact finite linear algebra in the
\path{verification/evidence/relative_cubic_depth.json} bundle, generated
by \path{relative_cubic_depth.py} and checked by its independent replay.
The fingerprint in Appendix G fixes all three files and their exact inputs.

\begin{proof}
Exact invariant and commutator calculations give
$\dim R=\dim M_G=3$, the two ranks in (F.1), and $N\pi=0$.
Thus $\operatorname{im}\pi\subseteq\ker N$; both spaces have dimension
two, so they are equal.  The computed image of $N$ is the displayed
kernel of $\pi$, proving (F.2).

The invariant symmetric-form pencil on $W$ has a unique rank-one member
$\ell^2$, where $\ell$ transforms by $\chi$.  Contracting a
$\chi$-twisted cubic by $\ell$ cancels the two characters and lands in
the invariant quadratic pencil.  Exact contraction gives (F.3), whose
kernel proves the second assertion.
\end{proof}

The kernel line is intrinsic rather than an artifact of the displayed
basis.  Among the $133$ projective lines of $R$, it is the unique line
whose first flattening has rank $9$; the remaining census is one line
of rank $1$ and $131$ lines of rank $10$.

\begin{remark}[Why the two canonical planes do not glue]
\label{rem:tate-depth-boundary}
The three labelled cubic orbit classes project to $M_G$ with relation
$[2,9,1]$, whereas the three labelled depth profiles have relation
$[2,8,1]$.  There is therefore no label-preserving linear map carrying
these source classes to the depth rays.  The integral transfer makes
the obstruction structural.  In orbit-size order set
\[
 s=(1,1,1)^t,\qquad B=s(1,4,6).
\]
Then $B^2=11B$.  Divided transfer kills every integral balanced
relation $a$ with $(1,4,6)a=0$, but $Bs=11s$, so it fixes the depth
socle after division by $11$.  It separates the source relation from
the depth socle instead of identifying them.  The doubled-line and
residual-line flag of Remark~\ref{rem:cubic-target-flag} is internal to
the depth plane and supplies no missing map.
\end{remark}

\section{Verification architecture}
\label{sec:verification}

Table~\ref{tab:verification} separates the mathematical arguments from
their exact finite inputs.  Here ``certificate'' means a deterministic
calculation over a prime field, checked against a canonical JSON output
and a SHA-256 manifest; an independent replay uses a separate
implementation.  ``Lean'' means that the stated implication is
checked by the Lean kernel.  Neither term replaces the cited
identification of the classical groups or configurations.

\begin{table}[ht]
\centering
\footnotesize
\begin{tabular}{@{}>{\raggedright\arraybackslash}p{0.24\textwidth}
                    >{\raggedright\arraybackslash}p{0.32\textwidth}
                    >{\raggedright\arraybackslash}p{0.34\textwidth}@{}}
\toprule
Result & Proof in the text or classical input & Exact evidence \\
\midrule
Matching-secant quotient
& unique conic divisibility and the four-endpoint switch
& none \\
Balanced-orbit completeness
& projective--trade bridge, Faber tame subgroups, targeted linear
  detectors, root-defect injection, and Steinberg/Fischer tests
& structural Lean corroboration and a \(q=9\) replay; neither replaces the
  all-field detector proofs in Appendix~\ref{app:modular-detectors} \\
Ranks and the middle Fischer layer
& affine cocycle, irreducible Fischer modules; Edge--Dye configurations
& quotient matrices as cross-checks; all edge-selected alternating cycles
  and independent Dickson replay \\
Balanced sheets and cubic orientation
& Propositions~\ref{prop:radical-hadamard} and
  \ref{prop:modular-sheet-mechanism};
  Lemma~\ref{lem:hyperplane-square}; maximal-isotropic quotient duality
& finite moments and ranks as cross-checks; kernel-checked
  hyperplane-square implication; Gorenstein replay \\
Fixed-line ambiguity and Chow rigidity
& invariant-space dimensions, subgroup normalizers, unique block systems,
  and radial translation
& structural Lean gate for the radial-translation implications; the
  coalescence-point square rank is outside the theorem \\
Six-profile reconstruction
& subgroup marks, orbit--stabilizer, and the three-ray lemmas
& profile rows and replay; decorated-parent bijection \\
Modular depth and the Tate plane
& projectivity, fixed-point exactness, and standard modular-module input
& invariant, coinvariant, and depth calculations with replay \\
Arithmetic splitting and gluing
& displayed root calculations; cited subgroup identifications
& Lean gate and normalized certificate/replay \\
\bottomrule
\end{tabular}
\caption{Proof modes for the principal results.}
\label{tab:verification}
\end{table}

The claim-by-claim map is
\path{verification/trust_manifest.json}; the exact statement
identity is \path{verification/statement_identity.json}.  From the root
of the archival source, the aggregate entry point is
\begin{verbatim}
python3 verification/verify_release.py
\end{verbatim}
It checks the statement identity and all admitted checksum manifests, runs
the admitted primary checks and independent replays, elaborates the
Paper-II-structural, arithmetic-gluing, Hilbert-symmetry, and
hyperplane-square Lean gates, and rebuilds this manuscript with the pinned
manuscript Nix environment and build checker.

The reproducible review source is pinned by
\begin{center}
\small\path{verification/evidence_fingerprint.json}.
\end{center}
Its SHA-256 digest is
\begin{center}
\small\texttt{0a5fc4a720769e6ae1251cff263eceb63a7824}\\[-2pt]
\texttt{8af2c85bd8862cb56734b93c96}.
\end{center}
That fingerprint fixes a normalized manuscript, statement contents,
extractor, paper Nix inputs, manuscript checker, adversarial boundary checker,
verification documentation, the admitted evidence manifests, command vectors,
Lean gates and their full project-owned import closure, runner, trust manifest, and
environment: Python 3.13.12, Lean
4.32.0-rc1, and Mathlib revision
\begin{center}
\small\path{571b8a8e54219b4d393f75f4b8653fac08197fcc}.
\end{center}
A successful replay ends with
\begin{center}
\texttt{clebsch factorization release: CHECK OK}.
\end{center}
The normalization replaces only the displayed fingerprint digest and
the statement identity's derived full-source hash, avoiding a circular
self-hash.  This is an explicit review-source allowlist, not a claim to
hash every file in the repository or every transitive tool dependency.
The fingerprint identifies the source closure; it is not an immutable
public locator.

The Paper-II structural gate
\begin{center}
\small\texttt{RelativeConicArcs.Gates.}\\[-2pt]
\texttt{ClebschPaperIIStructural}
\end{center}
follows the proof spine from the projective pullback through the concrete
finite-root and outer-defect abstractions, the divided-power detectors, and
the affine contraction, to the rank-three endpoint.  The all-field
representation and subgroup assertions in
Appendix~\ref{app:modular-detectors} remain human proofs.  Its twenty-nine
audited terminals include the derived fused
\(q=5\) exclusion and point-stabilizer uniqueness certificates for the
\(B_3/\F_7\) and \(H_3/\F_{11}\) matchings.  The \(H_3\) uniqueness check
uses twelve explicit point-stabilizer witnesses per sheet, not a matching
census.  The same gate imports
\begin{center}
\small\texttt{ClebschFixedLineRadialTranslation}.
\end{center}
The exact terminals are
\begin{center}\small
\texttt{annihilates\_hadamardSquare\_iff\_}\\[-2pt]
\texttt{eq\_sheetSignLine\_of\_noncoalescent}
and
\path{card_nonmatchingNoncoalescentParameters}.
\end{center}
They check the table-free Radical--Hadamard inheritance and the exact
\(q-2\) count after removing the matching and coalescence parameters.
The inheritance is derived, not assumed: the evaluation space is presented as
the parameter-independent top space together with the constant function and
the radial level, the two radial levels at parameters with nonzero outer
constant span the same line, and the three Radical--Hadamard hypotheses are
therefore imposed only at the reference parameter.  The gate terminal
\path{evaluationSpace_eq_reference} records that step.  Identifying the
geometric evaluation spaces with this presentation remains a human input.
The invariant-space, stabilizer, block-system, and unique-Chow-point
arguments remain human and classical inputs.  Steinberg, Doty--Henke, and
modular Hermite reciprocity supply the cited representation-theoretic
interfaces; the printed tilting, root-defect, and contraction arguments
prove the detector consequences.  Faber supplies the abstract tame subgroup
types, while Dickson--Giudici enters only in the later exact-order
maximal-overgroup step.  These are classical inputs rather than consequences
of the Lean gate.

The largest formal component is the gate
\begin{center}
\small\texttt{RelativeConicArcs.Gates.}\\[-2pt]
\texttt{ClebschArithmeticGluing}.
\end{center}
Its twenty-three terminals check the displayed roots, matching involutions, projective
orders, determinant halves, the $A_3/B_3$ orbit calculations, and the
split--fused conclusion.  For $H_3$, Lean checks the normalized
certificate rows, intersections, determinant classes, the $11+11$
signature split, and the shape of the bounded generation words.
Evaluation and completeness of the $660$-element coset and word
certificates remain exact generator/replay obligations.  The classical
names $S_4,A_4,A_5$, and $\operatorname{PSL}_2(11)$ are cited inputs:
no group name is inferred from an order.

The separate structural gate
\begin{center}
\small\texttt{RelativeConicArcs.Gates.}\\[-2pt]
\texttt{ClebschHilbertSymmetry}
\end{center}
retains an independent check of the Hilbert arithmetic: a symmetric
finite Hilbert function of total length \(2q\) whose first values are
\(1,q-1,q-1\) has socle degree three and final value one.  The
Gorenstein proof does not use this implication; it obtains the perfect
degree-\(1\)-by-degree-\(2\) pairing directly from maximal isotropy.

The other computations are load-bearing only at the points shown in
Table~\ref{tab:verification}.  The matching-module bundle supplies the
three finite quotient ranks and cross-checks the moment data.  The
\(H_3\)-equivariant-rank bundle checks the Sylow-normalizer weights,
\(A_5\)-fixed dimensions, sheet membership, and the top witness.  The
small-field trade bundle exhausts \(q=9\); its larger-field rows are
cross-checks only.  The shared-radial bundle
reconstructs every edge-selected pair, its alternating cycle, and its
deepest trace; an independent replay starts from the Dickson recurrence.
The balanced-sheet bundle independently reconstructs the
permutation-module conclusions, while the Gorenstein bundle checks the
ideals, deletion ranks, socles, and resolutions.  The profile-incidence
bundle computes one row per double coset; the decorated-parent bundle
identifies a singleton matching with its parent; the
relative-cubic-depth bundle computes the modular and Tate rank patterns;
and the arithmetic-gluing gate checks the normalized splitting data.
All use exact prime-field arithmetic.  The balanced-orbit classification
uses no field census: its executable component checks the local formulas
after the human modular identifications.  The general divisibility theorem,
radical--Hadamard and hyperplane-square lemmas, modular-permutation
mechanism, three-ray cubic lemma, and mass-zero factorization are proved in
the text.

This verification surface proves only the three marked configurations
defined here.  It does not supply an all-characteristic construction or a
canonical map from the relative-cubic Tate plane to the depth plane.

\section*{AI assistance disclosure}

This project was developed with extensive assistance from OpenAI Codex and
Anthropic Claude. Under the author's direction, the systems assisted
throughout with proof exploration, literature research, symbolic and finite
computation, code and formal-proof development, verification, and manuscript
drafting and revision. The author checked the resulting arguments,
computations, code, and cited sources, reviewed and edited all AI-assisted
material, and assumes responsibility for all content.

\begin{thebibliography}{99}

\bibitem{Atlas1985}
J.~H.~Conway, R.~T.~Curtis, S.~P.~Norton, R.~A.~Parker, and
R.~A.~Wilson,
\newblock \emph{Atlas of Finite Groups},
\newblock Clarendon Press, Oxford, 1985.

\bibitem{Edge1956}
W.~L.~Edge,
\newblock Conics and orthogonal projectivities in a finite plane,
\newblock \emph{Canadian Journal of Mathematics} \textbf{8} (1956),
362--382.

\bibitem{Dye1991}
R.~H.~Dye,
\newblock Hexagons, conics, $A_5$ and $\mathrm{PSL}_2(K)$,
\newblock \emph{Journal of the London Mathematical Society} (2)
\textbf{44} (1991), 270--286.
\newblock doi:10.1112/jlms/s2-44.2.270.

\bibitem{CameronKorchmaros1993}
P.~J.~Cameron and G.~Korchm\'aros,
\newblock One-factorizations of complete graphs with a doubly transitive
automorphism group,
\newblock \emph{Bulletin of the London Mathematical Society} \textbf{25}
(1993), 1--6.
\newblock doi:10.1112/blms/25.1.1.

\bibitem{KorchmarosPaceSonnino2018}
G.~Korchm\'aros, N.~Pace, and A.~Sonnino,
\newblock One-factorisations of complete graphs arising from ovals in finite
planes,
\newblock \emph{Journal of Combinatorial Theory, Series A} \textbf{160}
(2018), 62--83.
\newblock doi:10.1016/j.jcta.2018.06.006.

\bibitem{Han2025}
H.~Han,
\newblock Symmetric factorisations of complete graphs,
\newblock \emph{Mathematical Foundations of Computing} \textbf{8}
(2025), no.~3, 372--378.
\newblock doi:10.3934/mfc.2023046.

\bibitem{Giudici2007}
M.~Giudici,
\newblock Maximal subgroups of almost simple groups with socle
$\mathrm{PSL}(2,q)$,
\newblock arXiv:math/0703685, 2007.

\bibitem{GiudiciLiPotocnikPraeger2006}
M.~Giudici, C.~H.~Li, P.~Poto\v{c}nik, and C.~E.~Praeger,
\newblock Homogeneous factorisations of graphs and digraphs,
\newblock \emph{European Journal of Combinatorics} \textbf{27} (2006),
11--37.
\newblock doi:10.1016/j.ejc.2004.08.003.

\bibitem{Faber2023}
X.~Faber,
\newblock Finite \(p\)-irregular subgroups of \(\mathrm{PGL}_2(k)\),
\newblock \emph{La Matematica} \textbf{2} (2023), 479--522;
arXiv:1112.1999.

\bibitem{DotyHenke2002}
S.~Doty and A.~Henke,
\newblock Decomposition of tensor products of modular irreducibles for
\(\mathrm{SL}_2\),
\newblock arXiv:math/0205186.

\bibitem{McDowellWildon2022}
E.~J.~McDowell and M.~Wildon,
\newblock Modular plethystic isomorphisms for two-dimensional linear groups,
\newblock \emph{J. Algebra} \textbf{602} (2022), 441--483.
\newblock doi:10.1016/j.jalgebra.2022.02.025.

\bibitem{Steinberg1963}
R.~Steinberg,
\newblock Representations of algebraic groups,
\newblock \emph{Nagoya Mathematical Journal} \textbf{22} (1963), 33--56.
\newblock doi:10.1017/S0027763000011016.

\bibitem{RodriguezPajaresEtAl2025}
G.~Rodr\'iguez-Pajares, D.~Ruano, and F.~Salizzoni,
\newblock A combinatorial description of when a self-associated set of
points fails to be arithmetically Gorenstein,
\newblock arXiv:2512.16766v1, 2025.

\bibitem{RuddRigidity2026}
T.~Rudd,
\newblock \emph{Reconstructing the Clebsch code and its golden orientation from its
deep-hole syndrome locus},
\newblock working paper, 2026.

\bibitem{BraunNebe2025}
T.~Braun and G.~Nebe,
\newblock Orthogonal representations of \(\operatorname{SL}_2(q)\) in
defining characteristic,
\newblock \emph{Beitr\"age zur Algebra und Geometrie} \textbf{66}(3)
(2025), 717--725.
\newblock doi:10.1007/s13366-024-00763-w.

\end{thebibliography}

\end{document}
